% SenseNova-U1 study zone.

\section*{SenseNova-U1 / NEO-Unify 源码学习专区}
\phantomsection
\addcontentsline{toc}{section}{SenseNova-U1 / NEO-Unify 源码学习专区}

\begin{conceptbox}
\textbf{学习目标：} 从 AIGC 工程与研究结合视角，读懂 SenseNova-U1 这类\key{统一多模态理解与生成}路线：它为什么重要、模型结构怎么分成 understanding / generation pathway、推理代码怎么组织，以及后续论文/训练代码公开后应该如何继续追。\\
\textbf{当前主线：} 先建立\key{技术路线位置 $\rightarrow$ 仓库地图 $\rightarrow$ 推理任务入口 $\rightarrow$ 模型参数与 pathway $\rightarrow$ 源码阅读顺序}，再逐步深入具体模块。
\end{conceptbox}

% ============================================================
% SenseNova-U1 / NEO-Unify 源码学习路线图
% ============================================================

\section{学习路线图：SenseNova-U1 / NEO-Unify 源码阅读}

\begin{conceptbox}
\textbf{一句话定位：} SenseNova-U1 代表的是当前统一多模态模型里很值得跟的一条路线：\key{用一个原生多模态架构统一理解、推理与生成，而不是把视觉编码器、LLM、扩散生成器通过 adapter 松散拼起来。}\\
\textbf{源码阅读目标：} 先不急着钻每个模块实现，而是先读懂\key{公开仓库里有哪些入口、四类推理任务如何调用同一个模型、understanding / generation pathway 如何分工、MoT 与 flow-matching 相关模块在代码里如何落位}。
\end{conceptbox}

\subsection*{资料收集与引用口径}

\begin{warnbox}
\textbf{当前资料状态：} 截至本笔记创建时，官方仓库 README 标注 arXiv / technical report 为 Coming，ToDo 中也写明训练代码和最终版权重/技术报告尚未完全公开。因此这一阶段的笔记只基于\key{官方仓库、官方中文 README、Hugging Face 模型集合、NEO-Unify 官方博客入口、公开推理示例与参数分解文档}建立学习地图，不把未公开论文细节当成既定事实。
\end{warnbox}

这一页主要综合以下公开资料：
\begin{itemize}[leftmargin=1.5em]
  \item OpenSenseNova 官方仓库：确认 SenseNova-U1 的项目定位、开源模型列表、公开代码范围和仓库目录。
  \item 官方中文 README：确认 NEO-Unify 的核心表述、模型规格、训练阶段描述、发布时间线与 ToDo 状态。
  \item \texttt{examples/README\_CN.md}：确认当前公开推理脚本覆盖 \texttt{t2i}、\texttt{editing}、\texttt{interleave}、\texttt{vqa} 四类任务。
  \item \texttt{docs/parameter\_breakdown\_CN.md}：确认 \texttt{SenseNova-U1-8B-MoT} 的参数拆分方式：understanding transformer、generation transformer 与 shared 文本输入输出层。
  \item Hugging Face 模型集合与 NEO-Unify 官方博客入口：作为模型权重、路线介绍和后续资料更新入口。
\end{itemize}

\subsection{一、为什么要单独建这个专区}

\begin{conceptbox}
\textbf{原因：} SenseNova-U1 不只是一个“文生图模型”或“VQA 模型”，而是一个\key{统一理解与生成}项目。它横跨模型架构、推理代码、图文交错生成、图像编辑、VQA、flow matching、MoT/MoE、prompt enhancement 和后续 RL 训练，因此更适合独立成“论文源码阅读”专区，而不是塞进某个单一主题目录。
\end{conceptbox}

从学习价值看，它至少连接了四条主线：
\begin{itemize}[leftmargin=1.5em]
  \item \textbf{模型架构主线：} 如何用 monolithic / native multimodal architecture 统一视觉理解与视觉生成。
  \item \textbf{生成建模主线：} 生成侧不再只是传统 diffusion pipeline，而是和语言/视觉 token、flow-matching head 等模块结合。
  \item \textbf{工程源码主线：} 公开仓库首先释放的是推理代码，因此最适合从 inference path 开始读。
  \item \textbf{后训练主线：} README 提到 SFT 模型经过理解预热、生成预训练、统一中期训练、统一 SFT，最终模型又经过一轮 T2I RL；这部分等技术报告或训练代码公开后再深入。
\end{itemize}

\subsection{二、当前公开信息里的核心口径}

\begin{summarybox}
\textbf{先记住这句话：} SenseNova-U1 想解决的不是“给 LLM 外挂一个图像生成器”，而是\key{让理解、推理和生成在一个原生多模态模型里发生}。官方对 NEO-Unify 的核心描述是：不依赖传统视觉编码器（VE）和变分自编码器（VAE）的模态翻译式拼接，而是端到端建模语言与视觉信息。
\end{summarybox}

需要暂时分清三层说法：
\begin{itemize}[leftmargin=1.5em]
  \item \textbf{官方 claim：} 统一多模态理解、推理与生成；从 modality integration 走向 true unification。
  \item \textbf{源码可验证部分：} 公开推理脚本、模型加载、任务入口、参数分组、generation / understanding pathway 的调用关系。
  \item \textbf{暂不展开部分：} 训练细节、完整技术报告、训练代码和最终版权重仍需等待官方进一步公开。
\end{itemize}

\subsection{三、仓库地图：先知道从哪里读}

当前仓库顶层结构可以先按功能理解：

\begin{center}
{\small
\begin{tabular}{p{3.2cm}p{9.3cm}}
\hline
\textbf{路径} & \textbf{源码阅读意义} \\
\hline
\texttt{README.md / README\_CN.md} & 项目定位、模型列表、公开状态、能力展示和后续 ToDo。 \\
\texttt{examples/} & 当前最重要的入口；包含文生图、图像编辑、图文交错生成、VQA 的参考推理脚本。 \\
\texttt{src/} & 模型与包代码主体；后续需要重点追模型类、processor、generation / understanding pathway、flow-matching 相关模块。 \\
\texttt{docs/} & 参数分解、prompt enhancement、base vs distill 等说明文档。 \\
\texttt{evaluation/} & 评测相关代码或配置；后续用于理解官方如何验证理解/生成能力。 \\
\texttt{scripts/} & 辅助脚本，例如参数统计脚本，是理解模型结构命名的好入口。 \\
\texttt{pyproject.toml / uv.lock} & Python 包依赖、项目安装与环境约束。 \\
\hline
\end{tabular}
}
\end{center}

\begin{intubox}
\textbf{源码阅读顺序不要反：} 第一轮不要直接冲进 \texttt{src/} 深处找所有 class。更稳的路线是：先从 \texttt{examples/} 的任务脚本看模型怎么被调用，再顺着调用链进入 \texttt{src/}，最后回头看参数分组和模块命名。
\end{intubox}

\subsection{四、第一轮源码阅读问题地图}

\begin{center}
{\small
\begin{tabular}{p{3.5cm}p{4.7cm}p{5.0cm}}
\hline
\textbf{核心问题} & \textbf{为什么重要} & \textbf{第一轮应该看哪里} \\
\hline
模型如何加载 & 所有任务都从这里开始，决定 processor、dtype、device、权重结构如何进入运行态 & \texttt{examples/*/inference.py}，再追到 \texttt{src/} 中的模型加载代码 \\
文生图怎么跑 & 这是生成侧主路径，能看到 prompt、CFG、timestep shift、num steps、输出分辨率等关键参数 & \texttt{examples/t2i/inference.py} \\
图像编辑怎么跑 & 这是图像条件生成路径，能看到输入图像 resize、image CFG、多图参考等工程逻辑 & \texttt{examples/editing/inference.py} \\
图文交错怎么跑 & 这是最能体现“统一理解与生成”的路径：一次响应里既有文本又有图像 & \texttt{examples/interleave/inference.py} \\
VQA 怎么跑 & 这是理解侧主路径，适合对照 generation pathway 看差异 & \texttt{examples/vqa/inference.py} \\
参数如何分组 & 帮助把名字里的 \texttt{understanding}、\texttt{generation}、\texttt{shared} 对上真实参数 & \texttt{docs/parameter\_breakdown\_CN.md} 与 \texttt{scripts/inspect\_model\_params.py} \\
\hline
\end{tabular}
}
\end{center}

\subsection{五、任务入口先验：四类推理脚本}

\subsubsection*{1. Text-to-Image：生成侧主路径}

\begin{itemize}[leftmargin=1.5em]
  \item \textbf{入口：} \texttt{examples/t2i/inference.py}
  \item \textbf{关键词：} \texttt{prompt}、\texttt{width}、\texttt{height}、\texttt{cfg\_scale}、\texttt{cfg\_norm}、\texttt{timestep\_shift}、\texttt{num\_steps}、\texttt{think}
  \item \textbf{学习重点：} 文本 prompt 如何进入模型，什么时候先生成 \texttt{<think>...</think>}，生成图像时哪些参数控制采样/去噪过程。
\end{itemize}

\subsubsection*{2. Image Editing：图像条件生成路径}

\begin{itemize}[leftmargin=1.5em]
  \item \textbf{入口：} \texttt{examples/editing/inference.py}
  \item \textbf{关键词：} \texttt{image}、\texttt{img\_cfg\_scale}、\texttt{input\_max\_pixels}、\texttt{target\_pixels}、\texttt{smart\_resize}
  \item \textbf{学习重点：} 输入图像如何预处理到模型可接受的尺寸，图像条件和文本条件如何共同影响生成。
\end{itemize}

\subsubsection*{3. Interleave：统一能力最值得看的路径}

\begin{itemize}[leftmargin=1.5em]
  \item \textbf{入口：} \texttt{examples/interleave/inference.py}
  \item \textbf{关键词：} \texttt{model.interleave\_gen}、\texttt{<think>}、文本输出、图像输出、JSONL 批量输入
  \item \textbf{学习重点：} 单次响应中如何交错地产生文本和图像；这是后面理解“native multimodal”最关键的源码路径。
\end{itemize}

\subsubsection*{4. VQA：理解侧主路径}

\begin{itemize}[leftmargin=1.5em]
  \item \textbf{入口：} \texttt{examples/vqa/inference.py}
  \item \textbf{关键词：} \texttt{image}、\texttt{question}、\texttt{max\_new\_tokens}、\texttt{do\_sample}、\texttt{temperature}、\texttt{top\_p}
  \item \textbf{学习重点：} 图像如何进入 understanding pathway，回答文本如何解码，和生成图像路径相比少了哪些 flow / denoising 逻辑。
\end{itemize}

\subsection{六、参数分解：先建立模型结构直觉}

\begin{summarybox}
\textbf{官方参数分解口径：} \texttt{SenseNova-U1-8B-MoT} 总参数约 17.552B；互斥分组包括 \texttt{generation\_transformer} 约 8.186B、\texttt{understanding\_transformer} 约 8.121B、\texttt{shared} 约 1.245B。由于 \texttt{embed\_tokens} 与 \texttt{lm\_head} 被两条 pathway 共用，所以从 forward pathway 视角看，理解路径约 9.366B，生成路径约 9.431B。
\end{summarybox}

这给源码阅读一个非常重要的直觉：
\begin{itemize}[leftmargin=1.5em]
  \item \textbf{understanding pathway：} 更接近图像理解 / VQA 路线，图像与文本进入理解侧模块，最后输出文本 token。
  \item \textbf{generation pathway：} 更接近图像生成 / 图文交错路线，会涉及生成侧视觉模块、flow-matching head、timestep / noise embedders 等。
  \item \textbf{shared：} 文本 token 的输入输出层是两条路径的公共部件，这也是为什么 \texttt{lm\_head} 不能简单理解为只属于“理解侧”。
\end{itemize}

\begin{intubox}
\textbf{读代码时的抓手：} 看到参数名或模块名里出现 \texttt{\_mot\_gen}、\texttt{fm\_modules}、\texttt{vision\_model\_mot\_gen}、\texttt{fm\_head}，优先联想到 generation pathway；看到不带生成后缀的视觉/语言模块，优先联想到 understanding pathway。
\end{intubox}

\subsection{七、建议后续拆成的笔记树}

\begin{itemize}[leftmargin=1.5em]
  \item \texttt{00\_learning\_roadmap.tex}：这页，先建立学习地图与公开资料口径。
  \item \texttt{02\_inference\_entrypoints.tex}：四个推理脚本的共同模式与差异。
  \item \texttt{03\_model\_loading\_and\_processor.tex}：模型加载、processor、token/image 输入格式。
  \item \texttt{04\_understanding\_pathway.tex}：VQA / understanding pathway 源码阅读。
  \item \texttt{05\_generation\_pathway.tex}：T2I / editing / flow-matching generation pathway 源码阅读。
  \item \texttt{06\_interleaved\_generation.tex}：图文交错生成，重点看 \texttt{interleave\_gen}。
  \item \texttt{07\_parameter\_and\_module\_map.tex}：参数命名、模块分组、MoT/MoE 相关结构。
  \item \texttt{08\_training\_and\_rl\_notes.tex}：等待官方 technical report / training code 后再补。
\end{itemize}

\subsection{八、阶段性学习目标}

\begin{summarybox}
\textbf{第一阶段不追求复现训练，只追求读懂推理代码。} 完成后你应该能做到：\\
1. 说清楚 SenseNova-U1 为什么是统一多模态路线，而不只是一个图像生成 pipeline；\\
2. 说清楚公开仓库里四类推理任务分别从哪里进入；\\
3. 根据参数分解解释 understanding pathway、generation pathway 和 shared 的区别；\\
4. 顺着一个样例脚本追到模型调用处，并知道下一步该读哪个模块；\\
5. 知道哪些内容目前公开资料还不足，不能过早下结论。
\end{summarybox}

% ============================================================
\subsection{参考资料}
% ============================================================

\begingroup
\small
\sloppy
\begin{itemize}[leftmargin=1.5em]
  \item OpenSenseNova，\emph{SenseNova-U1: Unifying Multimodal Understanding and Generation with NEO-Unify Architecture}，官方 GitHub 仓库。\\
  \url{https://github.com/OpenSenseNova/SenseNova-U1}
  \item OpenSenseNova，\emph{SenseNova-U1 中文 README}。\\
  \url{https://github.com/OpenSenseNova/SenseNova-U1/blob/main/README_CN.md}
  \item OpenSenseNova，\emph{SenseNova-U1 examples README：公开推理脚本说明}。\\
  \url{https://github.com/OpenSenseNova/SenseNova-U1/blob/main/examples/README_CN.md}
  \item OpenSenseNova，\emph{SenseNova-U1 模型参数分解}。\\
  \url{https://github.com/OpenSenseNova/SenseNova-U1/blob/main/docs/parameter_breakdown_CN.md}
  \item SenseNova，\emph{SenseNova-U1 Hugging Face 模型集合}。\\
  \url{https://huggingface.co/collections/sensenova/sensenova-u1}
  \item SenseNova，\emph{NEO-Unify 官方博客入口}。\\
  \url{https://huggingface.co/blog/sensenova/neo-unify}
\end{itemize}
\endgroup

% ============================================================
% SenseNova-U1 / NEO-Unify inference entrypoints
% ============================================================

\section{推理入口：先从 \texttt{examples/} 看模型怎么被调用}

\begin{conceptbox}
\textbf{这一节的目标：} 不急着深入大模型内部实现，先把官方公开的四个推理入口读清楚。\key{源码阅读第一步不是看所有 class，而是看外部任务如何把输入组织好，然后调用模型的哪个高层方法。}\\
\textbf{当前阅读基准：} 本节基于官方仓库 commit \texttt{2e339f94e992464d468794286af34ca759cb7ac0} 的 \texttt{examples/} 与 \texttt{src/sensenova\_u1/models/neo\_unify/modeling\_neo\_chat.py}。
\end{conceptbox}

\subsection{一、先给出最终地图}

\begin{summarybox}
\textbf{四个 examples 脚本本质上都是薄封装。} 它们共同完成三件事：\key{加载 tokenizer / model，处理任务输入，然后调用模型对象上的一个高层推理方法}。真正值得后续深入追的函数都集中在 \texttt{modeling\_neo\_chat.py}：\texttt{t2i\_generate}、\texttt{it2i\_generate}、\texttt{interleave\_gen} 和 \texttt{chat}。
\end{summarybox}

\begin{center}
{\small
\begin{tabular}{p{3.0cm}p{4.0cm}p{4.0cm}p{3.2cm}}
\hline
\textbf{任务} & \textbf{example 入口} & \textbf{模型高层方法} & \textbf{路径直觉} \\
\hline
文生图 & \texttt{examples/t2i/inference.py} & \texttt{model.t2i\_generate(...)} & 生成侧主路径 \\
图像编辑 & \texttt{examples/editing/inference.py} & \texttt{model.it2i\_generate(...)} & 图像条件生成路径 \\
图文交错 & \texttt{examples/interleave/inference.py} & \texttt{model.interleave\_gen(...)} & 统一能力展示路径 \\
VQA & \texttt{examples/vqa/inference.py} & \texttt{model.chat(...)} & 理解侧主路径 \\
\hline
\end{tabular}
}
\end{center}

\begin{intubox}
\textbf{最重要的阅读姿势：} \texttt{examples/} 不是模型核心，它是“任务适配层”。它告诉我们：什么输入会被传进模型、哪些参数是用户可控的、输出如何被保存。下一步真正要钻的是 \texttt{src/sensenova\_u1/models/neo\_unify/modeling\_neo\_chat.py} 里对应的四个方法。
\end{intubox}

\subsection{二、四个入口的共同骨架}

四个推理脚本看起来任务不同，但外层套路高度一致：

\begin{enumerate}[leftmargin=1.8em]
  \item \textbf{解析命令行参数：} 都支持 \texttt{--model\_path}、\texttt{--device}、\texttt{--dtype}、\texttt{--attn\_backend}、\texttt{--profile} 等通用参数。
  \item \textbf{设置 attention 后端：} 通过 \texttt{sensenova\_u1.set\_attn\_backend(args.attn\_backend)} 控制使用 \texttt{auto}、\texttt{flash} 或 \texttt{sdpa}。
  \item \textbf{加载配置、tokenizer 和模型：} 基本模式是 \texttt{AutoConfig.from\_pretrained}、\texttt{AutoTokenizer.from\_pretrained}、\texttt{AutoModel.from\_pretrained}，再 \texttt{to(device).eval()}。
  \item \textbf{做任务输入预处理：} 文生图主要处理 prompt 和输出尺寸；编辑处理输入图像、缩放和输出尺寸；interleave 处理文本、可选输入图像和系统提示词；VQA 处理图像 tensor 与问题文本。
  \item \textbf{调用模型高层方法：} 入口脚本不自己实现 attention / transformer / flow matching，而是调用模型暴露出来的任务级方法。
  \item \textbf{保存输出：} 图像任务保存 \texttt{png}，think mode 额外保存推理文本；VQA 保存答案；interleave 同时保存文本和生成图像。
\end{enumerate}

可以把它抽象成下面这条流程：

\begin{center}
\begin{adjustbox}{max width=\linewidth}
\begin{tikzpicture}[
  node distance=1.6cm,
  box/.style={draw, rounded corners, align=center, inner sep=6pt, font=\small},
  arr/.style={->, thick}
]
\node[box] (cli) {CLI / JSONL\\输入参数};
\node[box, right=of cli] (load) {加载 tokenizer\\和 AutoModel};
\node[box, right=of load] (prep) {任务输入\\预处理};
\node[box, right=of prep] (call) {调用模型\\高层方法};
\node[box, right=of call] (save) {保存文本\\或图像输出};
\draw[arr] (cli) -- (load);
\draw[arr] (load) -- (prep);
\draw[arr] (prep) -- (call);
\draw[arr] (call) -- (save);
\end{tikzpicture}
\end{adjustbox}
\end{center}

\subsection{三、Text-to-Image：\texttt{t2i\_generate} 是生成侧主入口}

\begin{conceptbox}
\textbf{入口定位：} \texttt{examples/t2i/inference.py} 是最纯粹的生成侧路径。它只需要文本 prompt，就能走到模型的图像生成逻辑。
\end{conceptbox}

\texttt{SenseNovaU1T2I.generate} 里最关键的一步是调用：

\begin{lstlisting}[language=Python]
out = self.model.t2i_generate(
    self.tokenizer,
    prompt,
    image_size=image_size,
    cfg_scale=cfg_scale,
    cfg_norm=cfg_norm,
    timestep_shift=timestep_shift,
    cfg_interval=cfg_interval,
    num_steps=num_steps,
    batch_size=batch_size,
    seed=seed,
    think_mode=think_mode,
)
\end{lstlisting}

这里先不要把参数想复杂，先建立直觉：

\begin{itemize}[leftmargin=1.5em]
  \item \textbf{\texttt{prompt}：} 生成图像的文本条件。
  \item \textbf{\texttt{image\_size}：} 输出图像分辨率。
  \item \textbf{\texttt{cfg\_scale} / \texttt{cfg\_norm} / \texttt{cfg\_interval}：} classifier-free guidance 相关控制，决定条件约束在采样过程中怎么起作用。
  \item \textbf{\texttt{timestep\_shift} / \texttt{num\_steps}：} 和生成侧去噪 / flow-matching 采样过程有关。
  \item \textbf{\texttt{think\_mode}：} 如果打开，模型先生成 \texttt{<think>...</think>}，再做图像生成；脚本会把 think 文本另存为 \texttt{.think.txt}。
\end{itemize}

\begin{intubox}
\textbf{后续追源码时的抓手：} 看到 \texttt{t2i\_generate}，就把它理解为“文本条件 \(\rightarrow\) 生成图像”的主函数。下一步要看它如何组织 prompt token、图像 token、采样 timestep、CFG，以及最终如何把模型输出还原成图像 tensor。
\end{intubox}

\subsection{四、Image Editing：\texttt{it2i\_generate} 多了图像条件}

\begin{conceptbox}
\textbf{入口定位：} \texttt{examples/editing/inference.py} 是 image-to-image / image-text-to-image 路径。和 T2I 相比，它不只是给 prompt，还要把输入图像作为条件传给模型。
\end{conceptbox}

它的高层调用是：

\begin{lstlisting}[language=Python]
output = self.model.it2i_generate(
    self.tokenizer,
    prompt,
    list(images),
    image_size=image_size,
    cfg_scale=cfg_scale,
    img_cfg_scale=img_cfg_scale,
    cfg_norm=cfg_norm,
    timestep_shift=timestep_shift,
    cfg_interval=cfg_interval,
    num_steps=num_steps,
    batch_size=batch_size,
    think_mode=think_mode,
    seed=seed,
)
\end{lstlisting}

这条路径比 T2I 多出的重点是：

\begin{itemize}[leftmargin=1.5em]
  \item \textbf{输入图像列表：} \texttt{images} 会被转成 \texttt{list(images)} 传入模型，说明编辑任务支持多图输入。
  \item \textbf{\texttt{input\_max\_pixels}：} 控制输入图像预处理预算，避免输入图像过大。
  \item \textbf{\texttt{target\_pixels}：} 当没有显式指定输出宽高时，脚本会按第一张输入图像的长宽比，归一化到目标像素规模。
  \item \textbf{\texttt{img\_cfg\_scale}：} 这是编辑路径里额外出现的图像条件 CFG 权重，用来区分“文本条件”和“图像条件”的引导力度。
  \item \textbf{\texttt{compare}：} 脚本还支持保存输入图和输出图的对比图，方便观察编辑效果。
\end{itemize}

\begin{summarybox}
\textbf{一句话理解 editing：} T2I 是“按文字生成一张图”，editing 是“拿已有图像作为条件，再按文字要求改出一张图”。所以它在同样的生成采样参数之外，多了输入图像读取、缩放、尺寸对齐和 \texttt{img\_cfg\_scale}。
\end{summarybox}

\subsection{五、Interleave：\texttt{interleave\_gen} 是统一能力最值得看的路径}

\begin{conceptbox}
\textbf{入口定位：} \texttt{examples/interleave/inference.py} 是最能体现 SenseNova-U1 路线价值的 example。它不是单纯回答文字，也不是单纯生成图片，而是允许一次响应里同时出现文本和生成图像。
\end{conceptbox}

关键调用是：

\begin{lstlisting}[language=Python]
text, image_tensors = self.model.interleave_gen(
    self.tokenizer,
    prompt,
    images=list(input_images),
    image_size=image_size,
    cfg_scale=cfg_scale,
    img_cfg_scale=img_cfg_scale,
    timestep_shift=timestep_shift,
    cfg_interval=cfg_interval,
    num_steps=num_steps,
    system_message=system_message,
    think_mode=think_mode,
    seed=seed,
)
\end{lstlisting}

这条路径有几个明显特点：

\begin{itemize}[leftmargin=1.5em]
  \item \textbf{输出是两类：} 返回 \texttt{text} 和 \texttt{image\_tensors}，脚本再把图像 tensor 转成 PIL 图像保存。
  \item \textbf{可以带输入图像：} \texttt{--image} 是可重复参数，说明 interleave 可以在已有图像基础上做图文混合推理。
  \item \textbf{默认有 system message：} 默认系统提示词明确约束了 think / no-think 协议，以及 \texttt{<image1>}、\texttt{<image2>} 这类图像引用标签。
  \item \textbf{默认开启 think mode：} \texttt{--think\_mode} 默认是 true，也可以通过 \texttt{--no-think\_mode} 关闭。
  \item \textbf{图像尺寸策略：} 如果有输入图像，输出尺寸会按第一张输入图像经过 \texttt{smart\_resize} 对齐；如果没有输入图像，则使用预设分辨率 bucket 或显式宽高。
\end{itemize}

\begin{intubox}
\textbf{为什么它最值得看：} 如果只看 T2I，你会觉得它像图像生成模型；只看 VQA，你会觉得它像视觉理解模型。但 \texttt{interleave\_gen} 同时处理文本、输入图像、生成图像和 think 协议，是后续理解“统一多模态”最关键的源码入口。
\end{intubox}

\subsection{六、VQA：\texttt{chat} 是理解侧主入口}

\begin{conceptbox}
\textbf{入口定位：} \texttt{examples/vqa/inference.py} 是 visual understanding / VQA 路径。它的输出是文本答案，不涉及图像生成采样参数。
\end{conceptbox}

VQA 的核心流程是先把图像转成模型需要的视觉输入：

\begin{lstlisting}[language=Python]
pixel_values, grid_hw = load_image_native(image)
pixel_values = pixel_values.to(self.device, dtype=self.model.dtype)
grid_hw = grid_hw.to(self.device)
\end{lstlisting}

然后调用：

\begin{lstlisting}[language=Python]
response, updated_history = self.model.chat(
    self.tokenizer,
    pixel_values,
    question,
    generation_config,
    history=history,
    return_history=True,
    grid_hw=grid_hw,
)
\end{lstlisting}

和生成图像路径相比，VQA 的参数明显更像语言模型解码：

\begin{itemize}[leftmargin=1.5em]
  \item \textbf{\texttt{max\_new\_tokens}：} 控制最多生成多少回答 token。
  \item \textbf{\texttt{do\_sample} / \texttt{temperature} / \texttt{top\_p} / \texttt{top\_k}：} 控制文本采样策略。
  \item \textbf{\texttt{repetition\_penalty}：} 控制重复惩罚。
  \item \textbf{\texttt{history}：} 支持对话历史，说明它走的是 chat 风格接口。
  \item \textbf{\texttt{grid\_hw}：} 保留图像 patch 网格信息，让模型知道视觉 token 的空间布局。
\end{itemize}

\begin{summarybox}
\textbf{一句话理解 VQA：} 这条路径的核心不是“采样一张图”，而是“把图像编码成视觉 token，再和问题一起走文本回答解码”。因此它没有 \texttt{cfg\_scale}、\texttt{timestep\_shift}、\texttt{num\_steps} 这些图像生成采样参数。
\end{summarybox}

\subsection{七、把四条路径放在一起比较}

\begin{center}
{\small
\begin{tabular}{p{2.6cm}p{3.0cm}p{3.4cm}p{3.8cm}p{2.4cm}}
\hline
\textbf{路径} & \textbf{输入} & \textbf{输出} & \textbf{关键参数} & \textbf{下一步追踪} \\
\hline
\texttt{t2i} & prompt & image & \texttt{cfg\_scale}, \texttt{cfg\_norm}, \texttt{timestep\_shift}, \texttt{num\_steps}, \texttt{think} & \texttt{t2i\_generate} \\
\texttt{editing} & prompt + images & edited image & \texttt{img\_cfg\_scale}, \texttt{input\_max\_pixels}, \texttt{target\_pixels}, \texttt{think} & \texttt{it2i\_generate} \\
\texttt{interleave} & prompt + optional images & text + images & \texttt{system\_message}, \texttt{think\_mode}, \texttt{img\_cfg\_scale}, \texttt{num\_steps} & \texttt{interleave\_gen} \\
\texttt{vqa} & image + question & answer text & \texttt{max\_new\_tokens}, \texttt{do\_sample}, \texttt{temperature}, \texttt{top\_p} & \texttt{chat} \\
\hline
\end{tabular}
}
\end{center}

\subsection{八、源码阅读顺序建议}

\begin{enumerate}[leftmargin=1.8em]
  \item \textbf{先读 \texttt{examples/t2i/inference.py}：} 它最干净，能先建立“加载模型 \(\rightarrow\) 调 \texttt{t2i\_generate} \(\rightarrow\) 保存图像”的基本骨架。
  \item \textbf{再读 \texttt{examples/vqa/inference.py}：} 它和 T2I 形成强对照：一个生成图像，一个生成文本答案。
  \item \textbf{接着读 \texttt{examples/editing/inference.py}：} 理解图像条件如何参与生成，尤其是输入 resize、输出尺寸和 \texttt{img\_cfg\_scale}。
  \item \textbf{最后读 \texttt{examples/interleave/inference.py}：} 这里最复杂，也最能体现统一多模态能力，适合在 T2I / VQA / Editing 都有直觉之后再看。
  \item \textbf{进入 \texttt{modeling\_neo\_chat.py}：} 分别追 \texttt{t2i\_generate}、\texttt{chat}、\texttt{it2i\_generate}、\texttt{interleave\_gen} 的内部逻辑。
\end{enumerate}

\begin{warnbox}
\textbf{暂时不要做的事：} 不要一上来就把 \texttt{modeling\_neo\_chat.py} 从头到尾硬读。这个文件很大，直接读会迷路。更好的方法是从四个 example 的调用点切进去，每次只追一个任务路径。
\end{warnbox}

\subsection{九、本节小结}

\begin{summarybox}
SenseNova-U1 的公开推理入口可以先理解成四条任务线：\key{T2I 调 \texttt{t2i\_generate}，Editing 调 \texttt{it2i\_generate}，Interleave 调 \texttt{interleave\_gen}，VQA 调 \texttt{chat}}。\texttt{examples/} 的价值不是展示模型内部细节，而是告诉我们每个任务如何准备输入、暴露哪些推理参数、最终进入哪个模型高层方法。
\end{summarybox}

% ============================================================
\subsection{参考资料}
% ============================================================

\begingroup
\small
\sloppy
\begin{itemize}[leftmargin=1.5em]
  \item OpenSenseNova，\emph{SenseNova-U1 examples README：公开推理脚本说明}。\\
  \url{https://github.com/OpenSenseNova/SenseNova-U1/blob/main/examples/README_CN.md}
  \item OpenSenseNova，\emph{Text-to-Image inference script}。\\
  \url{https://github.com/OpenSenseNova/SenseNova-U1/blob/main/examples/t2i/inference.py}
  \item OpenSenseNova，\emph{Image Editing inference script}。\\
  \url{https://github.com/OpenSenseNova/SenseNova-U1/blob/main/examples/editing/inference.py}
  \item OpenSenseNova，\emph{Interleaved text-image inference script}。\\
  \url{https://github.com/OpenSenseNova/SenseNova-U1/blob/main/examples/interleave/inference.py}
  \item OpenSenseNova，\emph{Visual understanding / VQA inference script}。\\
  \url{https://github.com/OpenSenseNova/SenseNova-U1/blob/main/examples/vqa/inference.py}
  \item OpenSenseNova，\emph{NEO-Unify model implementation：\texttt{modeling\_neo\_chat.py}}。\\
  \url{https://github.com/OpenSenseNova/SenseNova-U1/blob/main/src/sensenova_u1/models/neo_unify/modeling_neo_chat.py}
\end{itemize}
\endgroup

% ============================================================
% SenseNova-U1 / NEO-Unify model loading and processor notes
% ============================================================

\section{模型加载与输入格式：\texttt{AutoModel}、tokenizer 与图像 patch}

\begin{conceptbox}
\textbf{这一节的目标：} 读懂 SenseNova-U1 推理代码里“模型是怎么加载起来的、文本和图像是怎么喂进去的”。这里先给结论：\key{官方代码没有单独暴露一个传统意义上的 \texttt{AutoProcessor}；输入加工由 \texttt{AutoTokenizer}、\texttt{load\_image\_native}、聊天模板，以及模型方法内部的 \texttt{<image>} token 展开机制共同完成。}\\
\textbf{当前阅读基准：} 本节基于官方仓库 commit \texttt{2e339f94e992464d468794286af34ca759cb7ac0} 的 \texttt{examples/}、\texttt{src/sensenova\_u1/\_\_init\_\_.py}、\texttt{src/sensenova\_u1/models/neo\_unify/\_\_init\_\_.py}、\texttt{utils.py}、\texttt{configuration\_neo\_chat.py} 与 \texttt{modeling\_neo\_chat.py}。
\end{conceptbox}

\subsection{一、先给最终心智模型}

\begin{summarybox}
\textbf{SenseNova-U1 的输入链路可以概括成一句话：} \key{文本走 tokenizer，图像走 \texttt{load\_image\_native} 变成 patch 序列；文本 prompt 中的 \texttt{<image>} 占位符会被替换成 \texttt{<img>} + 多个 \texttt{<IMG\_CONTEXT>} + \texttt{</img>}，然后模型把 \texttt{<IMG\_CONTEXT>} 位置的文本 embedding 替换成视觉 embedding。}
\end{summarybox}

\begin{center}
\begin{adjustbox}{max width=\linewidth}
\begin{tikzpicture}[
  node distance=0.85cm and 0.75cm,
  box/.style={draw, rounded corners, align=center, inner sep=5pt, font=\small, text width=3.25cm},
  arr/.style={->, thick}
]
\node[box] (text) {prompt / question\\含可选 \texttt{<image>}};
\node[box, right=of text] (query) {conversation template\\生成 query};
\node[box, right=of query] (tokens) {\texttt{<image>} 展开成\\\texttt{<img>} + N 个\\\texttt{<IMG\_CONTEXT>} + \texttt{</img>}};
\node[box, below=of text] (image) {PIL image\\或图片路径};
\node[box, right=of image] (patch) {\texttt{load\_image\_native}\\得到 \texttt{pixel\_values}\\和 \texttt{grid\_hw}};
\node[box, right=of patch] (embed) {视觉特征替换\\\texttt{<IMG\_CONTEXT>} embedding};
\node[box, below=of embed] (model) {NEOChatModel\\理解或生成};
\draw[arr] (text) -- (query);
\draw[arr] (query) -- (tokens);
\draw[arr] (image) -- (patch);
\draw[arr] (patch) -- (embed);
\draw[arr] (tokens) -- (embed);
\draw[arr] (embed) -- (model);
\end{tikzpicture}
\end{adjustbox}
\end{center}

\begin{intubox}
\textbf{容易误解的点：} 很多多模态 Hugging Face 项目会有 \texttt{AutoProcessor.from\_pretrained(...)}，但 SenseNova-U1 当前公开代码不是这个组织方式。它把 processor 的职责拆散到了 tokenizer、图像工具函数和模型高层方法内部。
\end{intubox}

\subsection{二、模型加载：为什么 \texttt{AutoModel.from\_pretrained} 能直接用}

四个官方推理入口的加载代码几乎完全一致：

\begin{lstlisting}[language=Python]
config = AutoConfig.from_pretrained(model_path)
check_checkpoint_compatibility(config)
self.tokenizer = AutoTokenizer.from_pretrained(model_path)
self.model = AutoModel.from_pretrained(
    model_path,
    config=config,
    torch_dtype=dtype,
).to(device).eval()
\end{lstlisting}

这里关键不在这四行本身，而在 \texttt{import sensenova\_u1} 的副作用：包初始化时会调用内部注册函数，把 NEO-Unify 的 config / model 类型注册到 Transformers 的 \texttt{AutoConfig} 和 \texttt{AutoModel}。

\begin{lstlisting}[language=Python]
AutoConfig.register("neo_vision", NEOVisionConfig, exist_ok=True)
AutoConfig.register("neo_chat", NEOChatConfig, exist_ok=True)

AutoModel.register(NEOVisionConfig, NEOVisionModel, exist_ok=True)
AutoModel.register(NEOChatConfig, NEOChatModel, exist_ok=True)
\end{lstlisting}

\begin{summarybox}
\textbf{这意味着：} 只要先导入 \texttt{sensenova\_u1}，Transformers 看到 checkpoint 的 \texttt{model\_type = neo\_chat} 时，就知道应该实例化 \texttt{NEOChatConfig} 和 \texttt{NEOChatModel}。这也是官方 example 注释里说“不需要 \texttt{trust\_remote\_code=True}”的原因。
\end{summarybox}

\subsection{三、checkpoint 兼容检查：小但重要}

加载 config 后，example 会立刻调用：

\begin{lstlisting}[language=Python]
check_checkpoint_compatibility(config)
\end{lstlisting}

这个函数会读取 config 字段 \texttt{sensenova\_u1\_min\_version}。如果 checkpoint 声明需要更高版本的 \texttt{sensenova-u1} 包，而本地安装版本太旧，就直接抛出清晰错误。

\begin{intubox}
\textbf{源码阅读意义：} 这不是模型结构核心，但说明官方预期 checkpoint 和代码会一起演化。以后如果你换新权重、旧代码加载失败，第一时间应该检查 \texttt{sensenova\_u1\_min\_version} 和本地包版本。
\end{intubox}

\subsection{四、配置结构：\texttt{NEOChatConfig} 是组合配置}

\texttt{NEOChatConfig} 的 \texttt{model\_type} 是 \texttt{neo\_chat}，并且 \texttt{is\_composition = True}。它不是一个单纯 LLM config，而是组合了：

\begin{itemize}[leftmargin=1.5em]
  \item \textbf{\texttt{vision\_config}：} 对应 \texttt{NEOVisionConfig}，包含视觉侧 patch size 等配置。
  \item \textbf{\texttt{llm\_config}：} 对应 \texttt{NEOLLMConfig}，继承自 \texttt{Qwen3Config}，并额外加入 height / width rope 相关配置。
  \item \textbf{\texttt{downsample\_ratio}：} 默认 \texttt{0.5}，后面计算图像 token 数、空间索引和生成 token 网格时都会用到。
  \item \textbf{\texttt{template}：} 指定对话模板，后续通过 \texttt{get\_conv\_template(self.template)} 构造 prompt。
\end{itemize}

模型初始化时会保存：

\begin{lstlisting}[language=Python]
self.patch_size = config.vision_config.patch_size
self.template = config.template
self.downsample_ratio = config.downsample_ratio
self.conv_template = get_conv_template(self.template)
self.system_message = self.conv_template.system_message
\end{lstlisting}

这几个字段直接决定后面图像如何 patchify、\texttt{<image>} 要展开成多少个上下文 token，以及 query 使用什么聊天格式。

\subsection{五、文本侧：没有神秘 processor，本质是 conversation template + tokenizer}

T2I 路径用 \texttt{\_build\_t2i\_query} 构造 query：

\begin{lstlisting}[language=Python]
template = get_conv_template(self.template)
template.system_message = self.system_message if system_message is None else system_message
template.append_message(template.roles[0], prompt_text)
template.append_message(template.roles[1], None)
query = template.get_prompt()
\end{lstlisting}

VQA 路径也是同样思想：把历史对话和当前问题 append 到模板，再得到一个完整 prompt。区别是如果有图像但问题里没有写 \texttt{<image>}，\texttt{chat} 会自动在问题前面补：

\begin{lstlisting}[language=Python]
if history is None and pixel_values is not None and '<image>' not in question:
    question = '<image>\n' + question
\end{lstlisting}

随后调用 tokenizer：

\begin{lstlisting}[language=Python]
model_inputs = tokenizer(query, return_tensors='pt')
input_ids = model_inputs['input_ids'].to(self.device)
attention_mask = model_inputs['attention_mask'].to(self.device)
\end{lstlisting}

\begin{summarybox}
\textbf{文本侧 processor 的真实形态：} \key{conversation template 负责把 system / user / assistant 拼成模型熟悉的聊天格式；tokenizer 负责把 query 变成 token ids；模型内部再根据 token ids 计算时空索引。}
\end{summarybox}

\subsection{六、图像侧：\texttt{load\_image\_native} 返回两个核心对象}

VQA 入口的图像加载非常直接：

\begin{lstlisting}[language=Python]
pixel_values, grid_hw = load_image_native(image)
pixel_values = pixel_values.to(self.device, dtype=self.model.dtype)
grid_hw = grid_hw.to(self.device)
\end{lstlisting}

\texttt{load\_image\_native} 做了四件事：

\begin{enumerate}[leftmargin=1.8em]
  \item \textbf{读取并转 RGB：} 支持路径或 PIL 图像；如果是 RGBA，会先用白色背景合成成 RGB。
  \item \textbf{智能 resize：} \texttt{smart\_resize} 会让高宽都能被给定 factor 整除，并把总像素数限制在 \texttt{min\_pixels} 和 \texttt{max\_pixels} 之间。
  \item \textbf{归一化：} 使用 ImageNet mean/std 做 \texttt{ToTensor + Normalize}。
  \item \textbf{patchify：} 按 \texttt{patch\_size = 16} 默认逻辑，把图像展平成 patch 序列，并返回 \texttt{grid\_hw}。
\end{enumerate}

\texttt{preprocess\_pixel\_values} 的输出形状可以这样理解：

\begin{lstlisting}[language=Python]
# input:  pixel_values shape = [3, H, W]
# output: flatten_pixel_values shape = [grid_h * grid_w, 3 * patch_size ** 2]
# output: grid_hw shape = [[grid_h, grid_w]]
\end{lstlisting}

\begin{intubox}
\textbf{重点记住：} \texttt{pixel\_values} 不是常见的 \texttt{[B, 3, H, W]} 图像 batch，而是已经被 patchify 成二维 patch 序列；\texttt{grid\_hw} 则保留原始 patch 网格的高宽，后面用于恢复空间位置索引。
\end{intubox}

\subsection{七、\texttt{<image>} 如何变成可替换的视觉 token}

SenseNova-U1 用一个很关键的约定连接文本和图像：用户 prompt 里写的是 \texttt{<image>}，但模型真正看到的是：

\begin{lstlisting}[language=Python]
image_tokens = IMG_START_TOKEN + IMG_CONTEXT_TOKEN * num_patch_token + IMG_END_TOKEN
query = query.replace('<image>', image_tokens, 1)
\end{lstlisting}

其中默认 special token 是：

\begin{itemize}[leftmargin=1.5em]
  \item \textbf{\texttt{IMG\_START\_TOKEN}：} \texttt{<img>}
  \item \textbf{\texttt{IMG\_CONTEXT\_TOKEN}：} \texttt{<IMG\_CONTEXT>}
  \item \textbf{\texttt{IMG\_END\_TOKEN}：} \texttt{</img>}
\end{itemize}

\texttt{num\_patch\_token} 由图像网格和 \texttt{downsample\_ratio} 决定：

\begin{lstlisting}[language=Python]
num_patch_token = int(grid_hw[i, 0] * grid_hw[i, 1] * self.downsample_ratio**2)
\end{lstlisting}

在默认 \texttt{downsample\_ratio = 0.5} 时，\texttt{downsample\_ratio\^{}2 = 0.25}，也就是视觉 token 数相当于 patch 网格数量的四分之一。这和后续的 merge / downsample 机制一致。

\begin{summarybox}
\textbf{关键直觉：} \texttt{<image>} 只是用户态占位符；\texttt{<IMG\_CONTEXT>} 才是模型内部用于“放置视觉 embedding”的槽位。图像越大，\texttt{grid\_hw} 越大，需要的 \texttt{<IMG\_CONTEXT>} 槽位就越多。
\end{summarybox}

\subsection{八、视觉 embedding 如何塞回语言序列}

在 VQA 的 \texttt{generate} 里，如果传入了 \texttt{pixel\_values}，模型会先抽取视觉特征：

\begin{lstlisting}[language=Python]
vit_embeds = self.extract_feature(pixel_values, grid_hw=grid_hw)
\end{lstlisting}

然后找到文本 token 序列中所有 \texttt{<IMG\_CONTEXT>} 的位置，把这些位置原来的 token embedding 替换成视觉 embedding：

\begin{lstlisting}[language=Python]
input_embeds = self.language_model.get_input_embeddings()(input_ids)
selected = (input_ids == self.img_context_token_id)
assert selected.sum() != 0
input_embeds[selected] = vit_embeds.reshape(-1, C).to(input_embeds.device)
\end{lstlisting}

Image Editing 的 \texttt{\_build\_it2i\_inputs} 也做了同样的事情，只是它用于生成路径的 prefix 输入。

\begin{intubox}
\textbf{这就是“图文融合”的最小实现抓手：} 图像不是作为一个外部对象永远悬挂在旁边，而是经过视觉模型抽特征后，直接写入语言模型输入 embedding 序列中的 \texttt{<IMG\_CONTEXT>} 槽位。
\end{intubox}

\subsection{九、维度如何对齐：\texttt{dense\_embedding} 同时做降采样和投影}

这里有一个必须追清楚的问题：\key{ViT 提取出来的图像特征，为什么能直接替换语言模型的 token embedding？}因为语言模型输入 embedding 的最后一维固定是 \texttt{llm\_hidden\_size}，如果视觉特征最后一维不是同一个大小，下面这句赋值会直接形状不匹配：

\begin{lstlisting}[language=Python]
input_embeds[selected] = vit_embeds.reshape(-1, C).to(input_embeds.device)
\end{lstlisting}

SenseNova-U1 的对齐点在视觉侧的 \texttt{NEOVisionEmbeddings}。配置里显式保存了两个维度：

\begin{lstlisting}[language=Python]
hidden_size=1024
llm_hidden_size=2048
downsample_ratio=0.5
\end{lstlisting}

进入视觉 embedding 模块后，先用 \texttt{patch\_embedding} 把每个 \(16 \times 16\) 图像 patch 变成视觉隐空间向量：

\begin{lstlisting}[language=Python]
self.patch_embedding = nn.Conv2d(
    in_channels=config.num_channels,
    out_channels=self.embed_dim,
    kernel_size=self.patch_size,
    stride=self.patch_size,
)
\end{lstlisting}

此时每个 patch 的通道维度是视觉侧的 \texttt{hidden\_size}。随后真正完成“对齐到 LLM”的是 \texttt{dense\_embedding}：

\begin{lstlisting}[language=Python]
self.dense_embedding = nn.Conv2d(
    in_channels=self.embed_dim,
    out_channels=self.llm_embed_dim,
    kernel_size=self.downsample_factor,
    stride=self.downsample_factor,
)
\end{lstlisting}

它一口气做了两件事：

\begin{enumerate}[leftmargin=1.8em]
  \item \textbf{空间降采样：} \texttt{downsample\_ratio = 0.5}，所以 \texttt{downsample\_factor = 2}。一个 \(2 \times 2\) 的 patch 小块会合并成一个视觉 token，token 数量变成原来的四分之一。
  \item \textbf{通道投影：} \texttt{out\_channels = self.llm\_embed\_dim}，所以输出 token 的最后一维直接变成语言模型 hidden size。
\end{enumerate}

源码里还有两个 assert 把这个约束钉死：

\begin{lstlisting}[language=Python]
assert embeddings.shape[0] == int(patch_embeds.shape[0] / self.downsample_factor**2)
return embeddings
\end{lstlisting}

因此，图像特征能塞进 \texttt{<IMG\_CONTEXT>} 槽位，不是靠“运行时自动适配”，而是靠视觉模型末端提前产出和 LLM 输入 embedding 同维度的视觉 token。

\begin{summarybox}
\textbf{可以把它想成“坑和砖”的匹配：} \texttt{<IMG\_CONTEXT>} 在文本序列里提前挖好了若干个坑；\texttt{dense\_embedding} 负责把原始 patch 砖块压缩成同样数量、同样宽度的砖。坑的数量由 \texttt{grid\_hw * downsample\_ratio\^{}2} 决定；砖的宽度由 \texttt{llm\_hidden\_size} 决定。
\end{summarybox}

举一个小例子。假设原图被切成 \(4 \times 4 = 16\) 个 patch，\texttt{downsample\_ratio = 0.5}，则 \texttt{downsample\_factor = 2}：

\begin{lstlisting}[language=Python]
# before dense_embedding
patch_embeds: [16, hidden_size]

# after 2x2 merge + channel projection
vit_embeds: [4, llm_hidden_size]

# text side
<img> <IMG_CONTEXT> <IMG_CONTEXT> <IMG_CONTEXT> <IMG_CONTEXT> </img>
\end{lstlisting}

最后的结果是：4 个 \texttt{<IMG\_CONTEXT>} 槽位，对应 4 个 \texttt{vit\_embeds}；每个 \texttt{vit\_embeds} 的最后一维都等于语言模型输入 embedding 的最后一维。

\subsection{十、位置编码分两层：视觉侧 2D RoPE，LLM 侧 \texttt{t/h/w} 三轴 RoPE}

位置编码也要分清楚两件事：\key{视觉编码器内部先给 patch 注入二维空间感；进入 LLM 后，还要在统一图文序列里继续保留文本顺序和图像二维坐标。}

\subsubsection{视觉编码器内部：patch 先做 2D RoPE}

图像 patch 天然是二维网格。如果直接把它们拉平成一维序列，模型会丢掉“谁在左上、谁在右下”的结构。\texttt{NEOVisionEmbeddings} 里先根据 \texttt{grid\_hw} 生成每个 patch 的 \(x, y\) 坐标：

\begin{lstlisting}[language=Python]
abs_pos_x, abs_pos_y = build_abs_positions_from_grid_hw(grid_hw, device=patch_embeds.device)
\end{lstlisting}

然后把 embedding 最后一维拆成两半：一半按 \(x\) 坐标旋转，一半按 \(y\) 坐标旋转：

\begin{lstlisting}[language=Python]
x_part_1 = x[..., :dim_half]
x_part_2 = x[..., dim_half:]
rotated_part_1 = apply_rotary_emb_1d(x_part_1, cos_cached_x, sin_cached_x, abs_positions_x)
rotated_part_2 = apply_rotary_emb_1d(x_part_2, cos_cached_y, sin_cached_y, abs_positions_y)
return torch.cat((rotated_part_1, rotated_part_2), dim=-1)
\end{lstlisting}

通俗讲：一个 patch 不只是“第 7 个 patch”，而是“第 \(h\) 行、第 \(w\) 列的 patch”。二维 RoPE 就是在视觉特征里提前写入这种行列关系。

\subsubsection{进入 LLM 后：统一序列使用 \texttt{t/h/w} 三套索引}

视觉 token 替换进 \texttt{<IMG\_CONTEXT>} 后，模型面对的是一条混合序列：

\begin{lstlisting}[language=Python]
[文本 token, 文本 token, <IMG_CONTEXT>, <IMG_CONTEXT>, 文本 token]
\end{lstlisting}

这时只用普通一维位置还不够。SenseNova-U1 在语言模型 attention 里传入的是三行索引：

\begin{lstlisting}[language=Python]
indexes = torch.stack([t_indexes, h_indexes, w_indexes], dim=0)
\end{lstlisting}

其中：

\begin{itemize}[leftmargin=1.5em]
  \item \textbf{\texttt{t\_indexes}：} 控制文本/整体序列顺序；
  \item \textbf{\texttt{h\_indexes}：} 控制图像 token 的高度坐标；
  \item \textbf{\texttt{w\_indexes}：} 控制图像 token 的宽度坐标。
\end{itemize}

在 \texttt{modeling\_qwen3.py} 的 attention 里，query/key 会被拆成时间部分和空间部分：

\begin{lstlisting}[language=Python]
query_states_t, query_states_hw = query_states.chunk(2, dim=-1)
query_states_h, query_states_w = query_states_hw.chunk(2, dim=-1)
\end{lstlisting}

然后三部分分别做 RoPE：

\begin{lstlisting}[language=Python]
cos_t, sin_t = self.rotary_emb(hidden_states, indexes[0].unsqueeze(0))
cos_h, sin_h = self.rotary_emb_hw(hidden_states, indexes[1].unsqueeze(0))
cos_w, sin_w = self.rotary_emb_hw(hidden_states, indexes[2].unsqueeze(0))
\end{lstlisting}

最后再拼回完整的 query/key：

\begin{lstlisting}[language=Python]
query_states = torch.cat([query_states_t, query_states_h, query_states_w], dim=-1)
key_states = torch.cat([key_states_t, key_states_h, key_states_w], dim=-1)
\end{lstlisting}

所以这不是简单的“文本用文本位置，图片用图片位置”两套系统并排放着，而是在同一个 attention 里把每个 token 的位置拆成三条轴：\texttt{t} 管序列先后，\texttt{h/w} 管图像二维结构。

\subsubsection{一个通俗例子}

假设 prompt 被展开成：

\begin{lstlisting}[language=Python]
请 描述 这 张 图 ： <img> A B C D </img>
\end{lstlisting}

其中 \texttt{A/B/C/D} 是 4 个 \texttt{<IMG\_CONTEXT>}，对应一张 \(2 \times 2\) 图：

\begin{center}
\begin{tabular}{cc}
A: $(h=0,w=0)$ & B: $(h=0,w=1)$ \\
C: $(h=1,w=0)$ & D: $(h=1,w=1)$
\end{tabular}
\end{center}

文本侧的一维顺序告诉模型：

\begin{lstlisting}[language=Python]
请 -> 描述 -> 这 -> 张 -> 图 -> A -> B -> C -> D
\end{lstlisting}

图像侧的二维坐标告诉模型：

\begin{lstlisting}[language=Python]
A 在左上，B 在右上，C 在左下，D 在右下
\end{lstlisting}

把两者合起来，模型既知道“图片插在这句话的哪个位置”，也知道“每个视觉 token 在图片里的哪个位置”。这就是 \texttt{grid\_hw}、\texttt{t/h/w indexes} 和 RoPE 共同完成的事情。

\subsection{十一、\texttt{grid\_hw} 与三维索引：为什么模型知道图像空间位置}

模型不只需要知道“这里有视觉 token”，还需要知道这些 token 的二维空间位置。\texttt{get\_thw\_indexes} 会返回三行索引：

\begin{itemize}[leftmargin=1.5em]
  \item \textbf{\texttt{t\_indexes}：} 文本/序列维度上的位置。
  \item \textbf{\texttt{h\_indexes}：} 图像 token 的高度坐标。
  \item \textbf{\texttt{w\_indexes}：} 图像 token 的宽度坐标。
\end{itemize}

当 \texttt{grid\_hw} 存在，并且当前 token 是 \texttt{<IMG\_CONTEXT>} 时，模型会用 \texttt{build\_abs\_positions\_from\_grid\_hw} 给这些 token 填上空间坐标：

\begin{lstlisting}[language=Python]
selected = (input_ids == self.img_context_token_id)
abs_pos_w, abs_pos_h = build_abs_positions_from_grid_hw(
    grid_hw // int(1 / self.downsample_ratio),
    device=t_indexes.device,
)
h_indexes[selected] = abs_pos_h
w_indexes[selected] = abs_pos_w
\end{lstlisting}

这一步解释了为什么 \texttt{grid\_hw} 必须和 \texttt{pixel\_values} 一起传：\key{前者负责告诉模型视觉 token 的二维布局，后者负责提供视觉内容本身。}

\subsection{十二、不同任务里的输入格式差异}

\begin{center}
{\small
\begin{tabular}{p{2.6cm}p{3.2cm}p{3.8cm}p{4.4cm}}
\hline
\textbf{任务} & \textbf{文本输入} & \textbf{图像输入} & \textbf{进入模型前的关键加工} \\
\hline
VQA & question / chat query & 单图或多图 patch 序列 & 若 question 没有 \texttt{<image>}，自动补；再展开为 \texttt{<IMG\_CONTEXT>} 槽位 \\
T2I & prompt + generation system message & 无输入图像 & 构造 T2I query；后续根据输出分辨率建立生成图像 token 网格 \\
Editing & prompt + 一个或多个 \texttt{<image>} & 多张输入图 patch 序列 & 图像先 patchify；prompt 中 \texttt{<image>} 数量不足时自动补；视觉 embedding 写入槽位 \\
Interleave & prompt + optional images + system message & 可选多图 patch 序列 & 文本与输入图像混合；输出又可能包含文本和生成图像 \\
\hline
\end{tabular}
}
\end{center}

\begin{warnbox}
\textbf{不要把 \texttt{<image>} 和 \texttt{<img>} 混为一谈：} \texttt{<image>} 是 prompt 里的占位符；\texttt{<img>}、\texttt{<IMG\_CONTEXT>}、\texttt{</img>} 是模型真正 token 化前替换进去的 special token 序列。
\end{warnbox}

\subsection{十三、本节小结}

\begin{summarybox}
本节可以浓缩成五句话：\\
1. SenseNova-U1 通过本地注册机制接入 Transformers 的 \texttt{AutoConfig} / \texttt{AutoModel}，所以 example 里可以直接 \texttt{from\_pretrained}；\\
2. 当前公开代码没有单独的 processor 类，processor 职责被拆成 tokenizer、conversation template、\texttt{load\_image\_native} 和模型内部的 token 展开逻辑；\\
3. 图像输入的关键格式是 \texttt{pixel\_values + grid\_hw + <IMG\_CONTEXT>}，其中 \texttt{pixel\_values} 提供视觉内容，\texttt{grid\_hw} 提供空间布局，\texttt{<IMG\_CONTEXT>} 提供语言序列中的替换槽位；\\
4. 视觉特征能替换文本 embedding，是因为 \texttt{dense\_embedding} 同时把 patch 网格按 \texttt{downsample\_ratio} 降采样，并把通道数投影到 \texttt{llm\_hidden\_size}；\\
5. 位置编码分两层：视觉侧先对 patch 做 2D RoPE，LLM 侧再用 \texttt{t/h/w} 三轴 RoPE 同时表达文本顺序和图像二维坐标。
\end{summarybox}

% ============================================================
\subsection{参考资料}
% ============================================================

\begingroup
\small
\sloppy
\begin{itemize}[leftmargin=1.5em]
  \item OpenSenseNova，\emph{SenseNova-U1 官方 GitHub 仓库}。\\
  \url{https://github.com/OpenSenseNova/SenseNova-U1}
  \item OpenSenseNova，\emph{SenseNova-U1 examples README：公开推理脚本说明}。\\
  \url{https://github.com/OpenSenseNova/SenseNova-U1/blob/main/examples/README_CN.md}
  \item OpenSenseNova，\emph{T2I inference wrapper：\texttt{AutoConfig} / \texttt{AutoTokenizer} / \texttt{AutoModel} 加载方式}。\\
  \url{https://github.com/OpenSenseNova/SenseNova-U1/blob/main/examples/t2i/inference.py}
  \item OpenSenseNova，\emph{VQA inference wrapper：\texttt{load\_image\_native} 与 \texttt{model.chat}}。\\
  \url{https://github.com/OpenSenseNova/SenseNova-U1/blob/main/examples/vqa/inference.py}
  \item OpenSenseNova，\emph{Package init：\texttt{check\_checkpoint\_compatibility} 与模型注册入口}。\\
  \url{https://github.com/OpenSenseNova/SenseNova-U1/blob/main/src/sensenova_u1/__init__.py}
  \item OpenSenseNova，\emph{NEO-Unify AutoConfig / AutoModel registration}。\\
  \url{https://github.com/OpenSenseNova/SenseNova-U1/blob/main/src/sensenova_u1/models/neo_unify/__init__.py}
  \item OpenSenseNova，\emph{Image preprocessing utilities：\texttt{smart\_resize} 与 \texttt{load\_image\_native}}。\\
  \url{https://github.com/OpenSenseNova/SenseNova-U1/blob/main/src/sensenova_u1/models/neo_unify/utils.py}
  \item OpenSenseNova，\emph{NEOChatConfig：组合配置与 \texttt{downsample\_ratio}}。\\
  \url{https://github.com/OpenSenseNova/SenseNova-U1/blob/main/src/sensenova_u1/models/neo_unify/configuration_neo_chat.py}
  \item OpenSenseNova，\emph{NEOChatModel：\texttt{chat}、\texttt{t2i\_generate}、\texttt{it2i\_generate} 与输入 embedding 替换逻辑}。\\
  \url{https://github.com/OpenSenseNova/SenseNova-U1/blob/main/src/sensenova_u1/models/neo_unify/modeling_neo_chat.py}
  \item OpenSenseNova，\emph{NEOVisionModel：\texttt{patch\_embedding}、\texttt{dense\_embedding} 与视觉侧 2D RoPE 实现}。\\
  \url{https://github.com/OpenSenseNova/SenseNova-U1/blob/main/src/sensenova_u1/models/neo_unify/modeling_neo_vit.py}
  \item OpenSenseNova，\emph{Qwen3 attention modification：\texttt{t/h/w} 三轴 RoPE 与图文统一 attention}。\\
  \url{https://github.com/OpenSenseNova/SenseNova-U1/blob/main/src/sensenova_u1/models/neo_unify/modeling_qwen3.py}
  \item Alexey Dosovitskiy et al.，\emph{An Image is Worth 16x16 Words: Transformers for Image Recognition at Scale}。\\
  \url{https://arxiv.org/abs/2010.11929}
  \item Jianlin Su et al.，\emph{RoFormer: Enhanced Transformer with Rotary Position Embedding}。\\
  \url{https://arxiv.org/abs/2104.09864}
  \item Qwen Team，\emph{Qwen2-VL: Enhancing Vision-Language Model's Perception of the World at Any Resolution}。\\
  \url{https://arxiv.org/abs/2409.12191}
\end{itemize}
\endgroup

% ============================================================
% SenseNova-U1 / NEO-Unify visual understanding pathway notes
% ============================================================

\section{VQA / Understanding Pathway：从图片和问题到文本回答}

\begin{conceptbox}
\textbf{这一节的目标：} 沿着官方 VQA 推理脚本，把 SenseNova-U1 的 visual understanding 链路从外到内走一遍：\key{图片如何变成视觉 token，问题如何变成 chat query，\texttt{<IMG\_CONTEXT>} 如何被视觉 embedding 替换，\texttt{t/h/w} 位置索引如何进入 Qwen3 attention，最后如何自回归生成答案。}\\
\textbf{当前阅读基准：} 本节基于官方仓库 commit \texttt{2e339f94e992464d468794286af34ca759cb7ac0} 的 \texttt{examples/vqa/inference.py}、\texttt{modeling\_neo\_chat.py}、\texttt{modeling\_neo\_vit.py}、\texttt{modeling\_qwen3.py}、\texttt{utils.py} 与 \texttt{conversation.py}。
\end{conceptbox}

\subsection{一、先给最终调用栈}

\begin{summarybox}
\textbf{VQA 的主路径可以概括为：}
\[
\begin{aligned}
\texttt{answer()} &\rightarrow \texttt{load\_image\_native()} \rightarrow \texttt{model.chat()} \\
&\rightarrow \texttt{model.generate()} \rightarrow \texttt{extract\_feature()} \\
&\rightarrow \texttt{language\_model.generate()} \rightarrow \texttt{batch\_decode()}
\end{aligned}
\]
真正的“理解”发生在两次融合里：第一，视觉 encoder 把图片 patch 编码成和 LLM hidden size 对齐的视觉 token；第二，这些视觉 token 被写入文本序列里的 \texttt{<IMG\_CONTEXT>} 槽位，然后统一进入 Qwen3 decoder attention。
\end{summarybox}

\begin{center}
\begin{adjustbox}{max width=\linewidth}
\begin{tikzpicture}[
  node distance=1.35cm,
  box/.style={draw, rounded corners, align=center, inner sep=5pt, font=\small},
  arr/.style={->, thick}
]
\node[box] (input) {image + question};
\node[box, right=1.6cm of input] (pre) {\texttt{load\_image\_native}\\\texttt{pixel\_values}, \texttt{grid\_hw}};
\node[box, right=1.6cm of pre] (chat) {\texttt{chat()}\\拼 query + 展开 \texttt{<image>}};
\node[box, below=1.05cm of chat] (vision) {\texttt{extract\_feature()}\\视觉 token};
\node[box, right=1.6cm of chat] (replace) {替换\\\texttt{<IMG\_CONTEXT>} embedding};
\node[box, right=1.6cm of replace] (llm) {\texttt{Qwen3ForCausalLM.generate}\\自回归解码};
\node[box, right=1.6cm of llm] (out) {decode\\answer};
\draw[arr] (input) -- (pre);
\draw[arr] (pre) -- (chat);
\draw[arr] (chat) -- (replace);
\draw[arr] (pre) |- (vision);
\draw[arr] (vision) -- (replace);
\draw[arr] (replace) -- (llm);
\draw[arr] (llm) -- (out);
\end{tikzpicture}
\end{adjustbox}
\end{center}

\subsection{二、VQA wrapper：它只是很薄的一层}

官方入口 \texttt{examples/vqa/inference.py} 中的 \texttt{SenseNovaU1VQA.answer()} 主要做三件事：

\begin{lstlisting}[language=Python]
pixel_values, grid_hw = load_image_native(image)
pixel_values = pixel_values.to(self.device, dtype=self.model.dtype)
grid_hw = grid_hw.to(self.device)

response, updated_history = self.model.chat(
    self.tokenizer,
    pixel_values,
    question,
    generation_config,
    history=history,
    return_history=True,
    grid_hw=grid_hw,
)
\end{lstlisting}

也就是说，VQA wrapper 本身没有复杂逻辑。它负责：

\begin{itemize}[leftmargin=1.5em]
  \item \textbf{图像预处理：} 调用 \texttt{load\_image\_native} 得到 \texttt{pixel\_values} 和 \texttt{grid\_hw}；
  \item \textbf{设备与精度：} 把图像 patch 转到模型所在 device，并转成模型 dtype；
  \item \textbf{转交模型：} 把 tokenizer、图像 patch、问题、生成参数一起交给 \texttt{model.chat()}。
\end{itemize}

\begin{intubox}
\textbf{源码阅读抓手：} 从 \texttt{answer()} 开始读时，不要在 wrapper 里找“模型理解”的秘密。这里更像工程胶水，真正的核心在 \texttt{NEOChatModel.chat()} 和 \texttt{NEOChatModel.generate()}。
\end{intubox}

\subsection{三、第一步：图片先变成 \texttt{pixel\_values + grid\_hw}}

上一节已经讲过 \texttt{load\_image\_native} 的输入格式，这里从 VQA pathway 角度再强调一次：

\begin{lstlisting}[language=Python]
# pixel_values: [sum(grid_h * grid_w), 3 * patch_size * patch_size]
# grid_hw:      [[grid_h, grid_w], ...]
\end{lstlisting}

\texttt{pixel\_values} 保存 patch 内容，\texttt{grid\_hw} 保存 patch 原来的二维布局。理解链路里这两个对象一路传到 \texttt{generate()}：

\begin{lstlisting}[language=Python]
generation_output = self.generate(
    pixel_values=pixel_values,
    input_ids=input_ids,
    grid_hw=grid_hw,
    attention_mask=attention_mask,
    **generation_config
)
\end{lstlisting}

\begin{warnbox}
\textbf{不要只看 \texttt{pixel\_values}：} 如果丢掉 \texttt{grid\_hw}，模型仍然可能知道每个 patch 的内容，但不知道这些 patch 在图片里的行列关系；后面的 \texttt{h\_indexes} 和 \texttt{w\_indexes} 也就无法正确构造。
\end{warnbox}

\subsection{四、第二步：\texttt{chat()} 负责把 VQA 问题变成模型 query}

\texttt{chat()} 的第一处关键逻辑是自动补图像占位符：

\begin{lstlisting}[language=Python]
if history is None and pixel_values is not None and '<image>' not in question:
    question = '<image>\n' + question
\end{lstlisting}

这意味着单轮 VQA 时，用户即使只问“图里有什么？”，模型内部也会把问题改成类似：

\begin{lstlisting}[language=Python]
<image>
图里有什么？
\end{lstlisting}

然后它会设置图像 special token 的 id：

\begin{lstlisting}[language=Python]
img_context_token_id = tokenizer.convert_tokens_to_ids(IMG_CONTEXT_TOKEN)
self.img_context_token_id = img_context_token_id
self.img_start_token_id = tokenizer.convert_tokens_to_ids(IMG_START_TOKEN)
\end{lstlisting}

接下来，\texttt{chat()} 使用 conversation template 拼 prompt：

\begin{lstlisting}[language=Python]
template = get_conv_template(self.template)
template.system_message = self.system_message

for (old_question, old_answer) in history:
    template.append_message(template.roles[0], old_question)
    template.append_message(template.roles[1], old_answer)
template.append_message(template.roles[0], question)
template.append_message(template.roles[1], None)
query = template.get_prompt()
\end{lstlisting}

\begin{summarybox}
\textbf{到这一步为止：} 输入还只是“文本 query + 图片 patch”。图片还没有真正进入 LLM；query 里仍然只是一个 \texttt{<image>} 占位符。
\end{summarybox}

\subsection{五、第三步：\texttt{<image>} 展开成视觉 embedding 的槽位}

\texttt{chat()} 会根据 \texttt{grid\_hw} 计算需要多少个 \texttt{<IMG\_CONTEXT>}：

\begin{lstlisting}[language=Python]
for i in range(grid_hw.shape[0]):
    num_patch_token = int(
        grid_hw[i, 0] * grid_hw[i, 1] * self.downsample_ratio**2
    )
    image_tokens = IMG_START_TOKEN + IMG_CONTEXT_TOKEN * num_patch_token + IMG_END_TOKEN
    query = query.replace('<image>', image_tokens, 1)
\end{lstlisting}

这里有两个对齐关系：

\begin{itemize}[leftmargin=1.5em]
  \item \textbf{数量对齐：} \texttt{num\_patch\_token} 必须等于视觉 encoder 最终输出的 token 数；
  \item \textbf{语义对齐：} \texttt{<IMG\_CONTEXT>} 是“待替换槽位”，之后会被真正的视觉 embedding 覆盖。
\end{itemize}

举一个小例子：如果 \texttt{grid\_hw = [4, 4]}，并且 \texttt{downsample\_ratio = 0.5}，那么：

\begin{lstlisting}[language=Python]
num_patch_token = 4 * 4 * 0.5**2 = 4
query = query.replace(
    '<image>',
    '<img><IMG_CONTEXT><IMG_CONTEXT><IMG_CONTEXT><IMG_CONTEXT></img>',
    1,
)
\end{lstlisting}

随后 tokenizer 把整个 query 变成 \texttt{input\_ids} 和 \texttt{attention\_mask}。

\subsection{六、第四步：\texttt{generate()} 构造三轴位置索引}

进入 \texttt{NEOChatModel.generate()} 后，第一件重要的事不是抽视觉特征，而是先根据 \texttt{input\_ids} 和 \texttt{grid\_hw} 构造位置索引：

\begin{lstlisting}[language=Python]
indexes = self.get_thw_indexes(input_ids[0], grid_hw)
\end{lstlisting}

\texttt{get\_thw\_indexes()} 返回三行：

\begin{lstlisting}[language=Python]
return torch.stack([t_indexes, h_indexes, w_indexes], dim=0)
\end{lstlisting}

它们分别表示：

\begin{itemize}[leftmargin=1.5em]
  \item \textbf{\texttt{t\_indexes}：} token 在文本/整体序列里的时间或块位置；
  \item \textbf{\texttt{h\_indexes}：} 图像 token 在图片中的高度坐标；
  \item \textbf{\texttt{w\_indexes}：} 图像 token 在图片中的宽度坐标。
\end{itemize}

源码里最关键的一段是：

\begin{lstlisting}[language=Python]
selected = (input_ids == self.img_context_token_id)
abs_pos_w, abs_pos_h = build_abs_positions_from_grid_hw(
    grid_hw // int(1 / self.downsample_ratio),
    device=t_indexes.device,
)
h_indexes[selected] = abs_pos_h.to(t_indexes.device, t_indexes.dtype)
w_indexes[selected] = abs_pos_w.to(t_indexes.device, t_indexes.dtype)
\end{lstlisting}

\begin{intubox}
\textbf{通俗理解：} \texttt{<IMG\_CONTEXT>} 先在文本序列里占坑；\texttt{get\_thw\_indexes()} 再给每个坑贴上“这是第几行、第几列的图像块”的坐标标签。
\end{intubox}

\subsection{七、一个容易忽略的细节：图像 token 是一个 block}

\texttt{t\_indexes} 不是简单的 \texttt{0,1,2,3,...}。它由下面这几行构造：

\begin{lstlisting}[language=Python]
img_start_shift = torch.cat([
    torch.zeros(1, dtype=torch.long).to(input_ids.device),
    (input_ids == self.img_start_token_id).long(),
], dim=0)[:-1]
not_img_token = (input_ids != self.img_context_token_id).long()
t_indexes = ((img_start_shift + not_img_token).cumsum(0) - 1)
\end{lstlisting}

效果是：普通文本 token 会推进 \texttt{t}，但连续的 \texttt{<IMG\_CONTEXT>} 不会每个都单独推进一个新时间步。图像 token 更像被放进同一个“图像块”。

这和 Qwen3 里的 block causal mask 呼应：

\begin{lstlisting}[language=Python]
def create_block_causal_mask(index):
    mask = (idx_j == idx_i) | (arange.unsqueeze(0) <= arange.unsqueeze(1))
\end{lstlisting}

它的含义是：

\begin{itemize}[leftmargin=1.5em]
  \item \textbf{不同时间块之间：} 仍然遵守自回归因果顺序，后面的 token 能看前面；
  \item \textbf{同一个时间块内部：} 如果 \texttt{idx\_j == idx\_i}，可以相互注意，不强行按 patch 顺序单向看。
\end{itemize}

\begin{summarybox}
\textbf{为什么这很重要：} 图像 patch 本来就是二维整体，不是自然语言那样严格左到右生成的序列。把图像 token 放进同一个 block，可以让同一张图里的视觉 token 更自由地互相交互，同时又不破坏文本生成的因果性。
\end{summarybox}

\subsection{八、第五步：\texttt{extract\_feature()} 抽出视觉 token}

如果传入了 \texttt{pixel\_values}，\texttt{generate()} 会抽取视觉特征：

\begin{lstlisting}[language=Python]
if visual_features is not None:
    vit_embeds = visual_features
else:
    vit_embeds = self.extract_feature(pixel_values, grid_hw=grid_hw)
\end{lstlisting}

\texttt{extract\_feature()} 对 VQA 来说走的是理解侧 vision model：

\begin{lstlisting}[language=Python]
return self.vision_model(
    pixel_values=pixel_values,
    output_hidden_states=False,
    return_dict=True,
    grid_hw=grid_hw,
).last_hidden_state
\end{lstlisting}

视觉模型内部的关键路径是：

\begin{lstlisting}[language=Python]
pixel_values = pixel_values.view(-1, 3, self.patch_size, self.patch_size)
patch_embeds = self.gelu(self.patch_embedding(pixel_values)).view(-1, self.embed_dim)
patch_embeds = self._apply_2d_rotary_pos_emb(patch_embeds, grid_hw)
patches_per_img = self.dense_embedding(patches_per_img.permute(0, 3, 1, 2))
embeddings = torch.cat(patches_list, dim=0)
\end{lstlisting}

这里可以拆成四步：

\begin{enumerate}[leftmargin=1.8em]
  \item \textbf{patch reshape：} 把 \texttt{load\_image\_native} 产生的二维 patch 向量恢复成 \(3 \times 16 \times 16\) 小图块；
  \item \textbf{patch embedding：} 用 \texttt{Conv2d} 把每个 patch 投到视觉 hidden size；
  \item \textbf{视觉侧 2D RoPE：} 根据 \texttt{grid\_hw} 给 patch 注入二维位置；
  \item \textbf{\texttt{dense\_embedding}：} 同时完成空间降采样和通道投影，输出 LLM hidden size 的视觉 token。
\end{enumerate}

\subsection{九、第六步：视觉 token 替换 \texttt{<IMG\_CONTEXT>} embedding}

视觉 token 抽出来以后，\texttt{generate()} 会先拿到文本 token 的 embedding：

\begin{lstlisting}[language=Python]
input_embeds = self.language_model.get_input_embeddings()(input_ids)
B, N, C = input_embeds.shape
input_embeds = input_embeds.reshape(B * N, C)
\end{lstlisting}

然后找到所有 \texttt{<IMG\_CONTEXT>} 的位置：

\begin{lstlisting}[language=Python]
input_ids = input_ids.reshape(B * N)
selected = (input_ids == self.img_context_token_id)
assert selected.sum() != 0
input_embeds[selected] = vit_embeds.reshape(-1, C).to(input_embeds.device)
\end{lstlisting}

最后恢复成 batch 序列形状：

\begin{lstlisting}[language=Python]
input_embeds = input_embeds.reshape(B, N, C)
\end{lstlisting}

\begin{summarybox}
\textbf{这就是 VQA 中图文融合的关键瞬间：} 从 LLM 视角看，输入仍然是一条普通 embedding 序列；但其中若干个位置已经不再是 \texttt{<IMG\_CONTEXT>} 的文本 embedding，而是真正的视觉 embedding。
\end{summarybox}

\subsection{十、第七步：统一序列进入 \texttt{language\_model.generate()}}

完成替换后，\texttt{NEOChatModel.generate()} 调用语言模型的自回归生成：

\begin{lstlisting}[language=Python]
outputs = self.language_model.generate(
    inputs_embeds=input_embeds,
    indexes=indexes,
    attention_mask=attention_mask,
    generation_config=generation_config,
    output_hidden_states=output_hidden_states,
    use_cache=True,
    **generate_kwargs,
)
\end{lstlisting}

这里有两个参数特别关键：

\begin{itemize}[leftmargin=1.5em]
  \item \textbf{\texttt{inputs\_embeds}：} 已经混入视觉 token 的统一 embedding 序列；
  \item \textbf{\texttt{indexes}：} 三轴位置索引，告诉 Qwen3 attention 每个 token 的 \texttt{t/h/w} 坐标。
\end{itemize}

在 \texttt{Qwen3Model.forward()} 里，如果传入的是 \texttt{inputs\_embeds} 而不是 \texttt{input\_ids}，它会根据 \texttt{indexes[0]} 创建 block causal mask：

\begin{lstlisting}[language=Python]
causal_mask_mapping = {
    "full_attention": create_block_causal_mask(indexes[0]),
}
self.current_index = indexes[0].max()
\end{lstlisting}

后续自回归生成新文本 token 时，如果走 \texttt{input\_ids} 分支，模型会把新 token 当作普通文本位置继续推进：

\begin{lstlisting}[language=Python]
self.current_index += 1
indexes = torch.LongTensor([[self.current_index], [0], [0]]).to(input_ids.device)
\end{lstlisting}

\begin{intubox}
\textbf{直觉：} 首轮 prefilling 阶段使用完整的图文混合序列；后续 decoding 阶段一次生成一个文本 token，它只需要新的 \texttt{t} 位置，\texttt{h/w} 默认为 0。
\end{intubox}

\subsection{十一、第八步：Qwen3 attention 怎样用 \texttt{t/h/w}}

在 \texttt{Qwen3Attention.forward\_und()} 里，query/key 的 head dim 会被拆成三部分：

\begin{lstlisting}[language=Python]
query_states_t, query_states_hw = query_states.chunk(2, dim=-1)
query_states_h, query_states_w = query_states_hw.chunk(2, dim=-1)

key_states_t, key_states_hw = key_states.chunk(2, dim=-1)
key_states_h, key_states_w = key_states_hw.chunk(2, dim=-1)
\end{lstlisting}

然后分别按三套索引做 RoPE：

\begin{lstlisting}[language=Python]
cos_t, sin_t = self.rotary_emb(hidden_states, indexes[0].unsqueeze(0))
query_states_t, key_states_t = apply_rotary_pos_emb(query_states_t, key_states_t, cos_t, sin_t)

cos_h, sin_h = self.rotary_emb_hw(hidden_states, indexes[1].unsqueeze(0))
query_states_h, key_states_h = apply_rotary_pos_emb(query_states_h, key_states_h, cos_h, sin_h)

cos_w, sin_w = self.rotary_emb_hw(hidden_states, indexes[2].unsqueeze(0))
query_states_w, key_states_w = apply_rotary_pos_emb(query_states_w, key_states_w, cos_w, sin_w)
\end{lstlisting}

最后再拼回完整 query/key：

\begin{lstlisting}[language=Python]
query_states = torch.cat([query_states_t, query_states_h, query_states_w], dim=-1)
key_states = torch.cat([key_states_t, key_states_h, key_states_w], dim=-1)
\end{lstlisting}

\begin{summarybox}
\textbf{这一步的意义：} LLM attention 不只是看到“文本 token + 图像 token”的混合序列，还能在 attention 的位置编码里区分文本顺序 \texttt{t}、图像高度 \texttt{h} 和图像宽度 \texttt{w}。
\end{summarybox}

\subsection{十二、第九步：生成答案并解码}

\texttt{language\_model.generate()} 输出 token id 后，\texttt{chat()} 做最后的 decode：

\begin{lstlisting}[language=Python]
response = tokenizer.batch_decode(generation_output, skip_special_tokens=True)[0]
response = response.split(template.sep.strip())[0].strip()
history.append((question, response))
\end{lstlisting}

这里的 \texttt{template.sep.strip()} 来自 conversation template。\texttt{chat()} 会把它作为 \texttt{eos\_token\_id}，并在 decode 后用它截断回答：

\begin{lstlisting}[language=Python]
eos_token_id = tokenizer.convert_tokens_to_ids(template.sep.strip())
generation_config['eos_token_id'] = eos_token_id
\end{lstlisting}

\begin{warnbox}
\textbf{多轮 VQA 的细节：} 单轮时 \texttt{chat()} 会自动给 question 前面补 \texttt{<image>}。多轮时，旧的 question / answer 会从 \texttt{history} 重新拼进模板；如果第一轮已经带了 \texttt{<image>}，历史里也会保留这个占位符。理解多轮行为时，要同时看 \texttt{history}、\texttt{grid\_hw} 和 prompt 里有几个 \texttt{<image>}。
\end{warnbox}

\subsection{十三、用一个极简例子串起来}

假设用户输入：

\begin{lstlisting}[language=Python]
image = cat.png
question = "这张图里有什么？"
\end{lstlisting}

如果图片经过预处理后得到 \texttt{grid\_hw = [4, 4]}，且 \texttt{downsample\_ratio = 0.5}，那么链路是：

\begin{enumerate}[leftmargin=1.8em]
  \item \textbf{wrapper：} \texttt{load\_image\_native(cat.png)} 得到 16 个原始 patch 和 \texttt{grid\_hw=[4,4]}；
  \item \textbf{chat：} 自动把问题变成 \texttt{<image>\textbackslash n这张图里有什么？}；
  \item \textbf{占位展开：} \texttt{<image>} 变成 4 个 \texttt{<IMG\_CONTEXT>} 槽位；
  \item \textbf{视觉编码：} \texttt{NEOVisionModel} 把 16 个 patch 经 \texttt{dense\_embedding} 变成 4 个 LLM hidden size 的视觉 token；
  \item \textbf{embedding 替换：} 4 个视觉 token 覆盖 4 个 \texttt{<IMG\_CONTEXT>} 的文本 embedding；
  \item \textbf{位置索引：} 这 4 个视觉 token 获得 \(2 \times 2\) 的 \texttt{h/w} 坐标，并作为一个图像 block 参与 attention；
  \item \textbf{自回归生成：} Qwen3 decoder 基于图文混合上下文生成回答，例如“图中有一只猫”。
\end{enumerate}

\subsection{十四、调试 VQA pathway 时最值得检查的断点}

\begin{center}
{\small
\begin{tabular}{p{3.4cm}p{5.0cm}p{6.0cm}}
\hline
\textbf{断点位置} & \textbf{应该检查什么} & \textbf{常见问题} \\
\hline
\texttt{load\_image\_native} 后 & \texttt{pixel\_values.shape}、\texttt{grid\_hw} & 图像过大/过小导致 token 数异常，或 device/dtype 未迁移 \\
\texttt{chat()} 展开后 & query 里是否还有足够的 \texttt{<IMG\_CONTEXT>} & prompt 中 \texttt{<image>} 数量和输入图片数量不匹配 \\
\texttt{get\_thw\_indexes()} 后 & \texttt{indexes.shape} 是否为 \texttt{[3, seq\_len]} & \texttt{grid\_hw} 缺失导致 \texttt{h/w} 全为 0 \\
\texttt{extract\_feature()} 后 & \texttt{vit\_embeds.shape[-1]} 是否等于 LLM hidden size & 视觉 token 维度无法替换文本 embedding \\
替换 \texttt{input\_embeds[selected]} 时 & \texttt{selected.sum()} 是否等于视觉 token 数 & \texttt{<IMG\_CONTEXT>} 槽位数量和视觉 token 数不一致 \\
\texttt{language\_model.generate()} 后 & 是否按 \texttt{eos\_token\_id} 停止 & eos token 配置或模板分隔符异常导致输出截断不对 \\
\hline
\end{tabular}
}
\end{center}

\subsection{十五、本节小结}

\begin{summarybox}
VQA / understanding pathway 可以浓缩成五句话：\\
1. VQA wrapper 只负责加载图像、迁移 device/dtype、组装 generation config，然后调用 \texttt{model.chat()}；\\
2. \texttt{chat()} 把问题变成 conversation query，并把 \texttt{<image>} 展开成 \texttt{<img>}、若干 \texttt{<IMG\_CONTEXT>} 和 \texttt{</img>}；\\
3. \texttt{extract\_feature()} 通过 \texttt{NEOVisionModel} 把图像 patch 编成和 LLM hidden size 对齐的视觉 token；\\
4. \texttt{generate()} 把视觉 token 写入 \texttt{<IMG\_CONTEXT>} 槽位，并把 \texttt{t/h/w} 三轴索引一起传给 Qwen3；\\
5. Qwen3 decoder 在 block causal mask 和三轴 RoPE 的帮助下，对统一图文 embedding 序列做自回归生成，最后由 tokenizer decode 成文本回答。
\end{summarybox}

% ============================================================
\subsection{参考资料}
% ============================================================

\begingroup
\small
\sloppy
\begin{itemize}[leftmargin=1.5em]
  \item OpenSenseNova，\emph{SenseNova-U1 官方 GitHub 仓库}。\\
  \url{https://github.com/OpenSenseNova/SenseNova-U1}
  \item OpenSenseNova，\emph{SenseNova-U1 README：统一多模态理解、推理与生成架构说明}。\\
  \url{https://github.com/OpenSenseNova/SenseNova-U1/blob/main/README.md}
  \item OpenSenseNova，\emph{VQA inference wrapper：\texttt{SenseNovaU1VQA.answer}}。\\
  \url{https://github.com/OpenSenseNova/SenseNova-U1/blob/main/examples/vqa/inference.py}
  \item OpenSenseNova，\emph{NEOChatModel：\texttt{chat}、\texttt{generate}、\texttt{extract\_feature} 与 \texttt{get\_thw\_indexes}}。\\
  \url{https://github.com/OpenSenseNova/SenseNova-U1/blob/main/src/sensenova_u1/models/neo_unify/modeling_neo_chat.py}
  \item OpenSenseNova，\emph{NEOVisionModel：patch embedding、2D RoPE 与 \texttt{dense\_embedding}}。\\
  \url{https://github.com/OpenSenseNova/SenseNova-U1/blob/main/src/sensenova_u1/models/neo_unify/modeling_neo_vit.py}
  \item OpenSenseNova，\emph{Qwen3 attention modification：block causal mask 与 \texttt{t/h/w} 三轴 RoPE}。\\
  \url{https://github.com/OpenSenseNova/SenseNova-U1/blob/main/src/sensenova_u1/models/neo_unify/modeling_qwen3.py}
  \item OpenSenseNova，\emph{Image preprocessing utilities：\texttt{load\_image\_native} 与 \texttt{smart\_resize}}。\\
  \url{https://github.com/OpenSenseNova/SenseNova-U1/blob/main/src/sensenova_u1/models/neo_unify/utils.py}
  \item OpenSenseNova，\emph{Conversation template utilities：\texttt{get\_conv\_template} 与 prompt 拼接规则}。\\
  \url{https://github.com/OpenSenseNova/SenseNova-U1/blob/main/src/sensenova_u1/models/neo_unify/conversation.py}
  \item Alexey Dosovitskiy et al.，\emph{An Image is Worth 16x16 Words: Transformers for Image Recognition at Scale}。\\
  \url{https://arxiv.org/abs/2010.11929}
  \item Jianlin Su et al.，\emph{RoFormer: Enhanced Transformer with Rotary Position Embedding}。\\
  \url{https://arxiv.org/abs/2104.09864}
  \item Qwen Team，\emph{Qwen2-VL: Enhancing Vision-Language Model's Perception of the World at Any Resolution}。\\
  \url{https://arxiv.org/abs/2409.12191}
\end{itemize}
\endgroup

% ============================================================
% SenseNova-U1 / NEO-Unify generation pathway notes
% ============================================================

\section{T2I / Editing / Flow-Matching Generation Pathway：按源码分块读}

\begin{conceptbox}
\textbf{这一节的目标：} 不再先按概念拆 T2I / editing / interleave，而是\key{按官方源码本身的文件与函数顺序}阅读 generation pathway：先看三个 example wrapper 如何把参数转给模型，再看 \texttt{modeling\_fm\_modules.py} 里生成头的定义，接着看 \texttt{modeling\_neo\_chat.py} 的 helper，最后逐段读 \texttt{interleave\_gen()}、\texttt{it2i\_generate()} 和 \texttt{t2i\_generate()}。\\
\textbf{当前阅读基准：} 本节基于官方仓库 commit \texttt{2e339f94e992464d468794286af34ca759cb7ac0} 的 \texttt{examples/t2i/inference.py}、\texttt{examples/editing/inference.py}、\texttt{examples/interleave/inference.py}、\texttt{modeling\_neo\_chat.py} 与 \texttt{modeling\_fm\_modules.py}。
\end{conceptbox}

\begin{summarybox}
\textbf{本节读法：} 每一块先标出源码位置，再解释这段代码做什么。不要先记抽象结论；先沿着源码看到：wrapper 调哪个模型方法，模型方法怎么建 query/cache，怎么初始化噪声图，怎么在 loop 里调用 \texttt{\_t2i\_predict\_v()}，最后怎么 \texttt{unpatchify} 回图像。
\end{summarybox}

\subsection{一、example wrapper：三个入口都只是薄转发}

\subsubsection{1. \texttt{examples/t2i/inference.py:64--122}}

T2I wrapper 的核心代码在 \texttt{SenseNovaU1T2I.generate()}。它不做生成逻辑，只把 CLI / JSONL 解析出的参数交给 \texttt{model.t2i\_generate()}：

\begin{lstlisting}[language=Python]
out = self.model.t2i_generate(
    self.tokenizer,
    prompt,
    image_size=image_size,
    cfg_scale=cfg_scale,
    cfg_norm=cfg_norm,
    timestep_shift=timestep_shift,
    cfg_interval=cfg_interval,
    num_steps=num_steps,
    batch_size=batch_size,
    seed=seed,
    think_mode=think_mode,
)
\end{lstlisting}

所以从源码阅读角度，\texttt{examples/t2i/inference.py} 只回答两个问题：一是输出尺寸、CFG、步数、seed 这些参数从哪里来；二是模型返回的 normalized tensor 如何经过 \texttt{\_to\_pil()} 保存成图片。真正的 T2I 生成逻辑在 \texttt{modeling\_neo\_chat.py:1619--1787}。

\subsubsection{2. \texttt{examples/editing/inference.py:171--230}}

Editing wrapper 的 \texttt{SenseNovaU1Editing.edit()} 同样只是转发，但比 T2I 多了输入图片列表和 \texttt{img\_cfg\_scale}：

\begin{lstlisting}[language=Python]
output = self.model.it2i_generate(
    self.tokenizer,
    prompt,
    list(images),
    image_size=image_size,
    cfg_scale=cfg_scale,
    img_cfg_scale=img_cfg_scale,
    cfg_norm=cfg_norm,
    timestep_shift=timestep_shift,
    cfg_interval=cfg_interval,
    num_steps=num_steps,
    batch_size=batch_size,
    think_mode=think_mode,
    seed=seed,
)
\end{lstlisting}

因此 editing 的源码阅读重点不是 wrapper，而是 \texttt{it2i\_generate()} 里输入图像如何先被转成 \texttt{pixel\_values/grid\_hw}，再写入 \texttt{<IMG\_CONTEXT>} 槽位，最后变成 prefix cache。

\subsubsection{3. \texttt{examples/interleave/inference.py:119--167}}

Interleave wrapper 调的是 \texttt{model.interleave\_gen()}：

\begin{lstlisting}[language=Python]
text, image_tensors = self.model.interleave_gen(
    self.tokenizer,
    prompt,
    images=list(input_images),
    image_size=image_size,
    cfg_scale=cfg_scale,
    img_cfg_scale=img_cfg_scale,
    timestep_shift=timestep_shift,
    cfg_interval=cfg_interval,
    num_steps=num_steps,
    system_message=system_message,
    think_mode=think_mode,
    seed=seed,
)
\end{lstlisting}

这一点很重要：interleave 不是先完整生成文本再额外调用图像生成器，而是在 \texttt{interleave\_gen()} 的自回归循环中，遇到 \texttt{<img>} token 才切入图像 flow loop。

\subsection{二、\texttt{modeling\_fm\_modules.py}：生成头和时间步 embedding}

这部分源码定义的是 generation pathway 会用到的“小模块”，它们不是完整 pipeline，只是在 \texttt{modeling\_neo\_chat.py} 中被组装进 \texttt{fm\_modules}。

\subsubsection{1. \texttt{TimestepEmbedder}}

\texttt{TimestepEmbedder} 把标量时间步 \(t\) 变成和 LLM hidden size 对齐的 embedding。后面在采样 loop 里会看到它直接加到当前图像 token embedding 上：

\begin{lstlisting}[language=Python]
timestep_embeddings = self.fm_modules['timestep_embedder'](t_expanded)
image_embeds = image_embeds + timestep_embeddings
\end{lstlisting}

也就是说，时间条件不是单独进一个 scheduler 对象，而是作为 token embedding 的一部分进入 Qwen3 decoder。

\subsubsection{2. \texttt{fm\_head}}

\texttt{fm\_head} 的职责是把 Qwen3 输出的图像 token hidden state 投回 RGB patch 维度。普通分支在 \texttt{NEOChatModel.\_\_init\_\_()} 中构造为：

\begin{lstlisting}[language=Python]
fm_head = nn.Sequential(
    nn.Linear(llm_hidden_size, 4096, bias=True),
    nn.GELU(),
    nn.Linear(4096, output_dim, bias=True),
)
\end{lstlisting}

其中 \texttt{output\_dim} 是一块 RGB patch 的像素维度：

\begin{lstlisting}[language=Python]
merge_size = int(1 / self.downsample_ratio)
output_dim = 3 * (patch_size * merge_size) ** 2
\end{lstlisting}

源码读到这里先记住一句：\key{Qwen3 输出 hidden state，\texttt{fm\_head} 把 hidden state 变成 RGB patch 预测。}

\subsection{三、\texttt{modeling\_neo\_chat.py:366--623}：generation helper 逐块读}

\subsubsection{1. \texttt{patchify()} / \texttt{unpatchify()}：源码行 366--395}

生成 loop 里的图像状态是 \texttt{image\_prediction}，形状是 \texttt{[B, 3, H, W]}。\texttt{patchify()} 把它切成 patch 序列：

\begin{lstlisting}[language=Python]
def patchify(self, images, patch_size, channel_first=False):
    h, w = images.shape[2] // patch_size, images.shape[3] // patch_size
    x = images.reshape(shape=(images.shape[0], 3, h, patch_size, w, patch_size))
    if channel_first:
        x = torch.einsum('nchpwq->nhwcpq', x)
    else:
        x = torch.einsum('nchpwq->nhwpqc', x)
    x = x.reshape(shape=(images.shape[0], h * w, patch_size**2 * 3))
    return x
\end{lstlisting}

这里有两个用法，后面会反复出现：

\begin{itemize}[leftmargin=1.5em]
  \item \textbf{作为 flow 状态：} \texttt{patchify(image\_prediction, patch\_size * merge\_size)} 得到要被 Euler 更新的 \texttt{z}；
  \item \textbf{作为 vision encoder 输入：} \texttt{patchify(image\_prediction, patch\_size, channel\_first=True)} 得到送入生成侧 vision encoder 的 patch。
\end{itemize}

\texttt{unpatchify()} 做反向操作，把更新后的 patch 序列拼回 \texttt{[B,3,H,W]} 图像：

\begin{lstlisting}[language=Python]
def unpatchify(sle, x, patch_size, h=None, w=None):
    if h is None or w is None:
        h = w = int(x.shape[1]**.5)
    else:
        h = h // patch_size
        w = w // patch_size
    x = x.reshape(shape=(x.shape[0], h, w, patch_size, patch_size, 3))
    x = torch.einsum('nhwpqc->nchpwq', x)
    images = x.reshape(shape=(x.shape[0], 3, h * patch_size, w * patch_size))
    return images
\end{lstlisting}

\subsubsection{2. Euler 与时间表：源码行 397--429}

\texttt{\_euler\_step()} 很短：

\begin{lstlisting}[language=Python]
def _euler_step(self, v_pred, z, t, t_next):
    z_next = z + (t_next - t) * v_pred
    return z_next
\end{lstlisting}

主循环里没有直接调用这个 helper，而是展开写成：

\begin{lstlisting}[language=Python]
z = z + (t_next - t) * v_pred
\end{lstlisting}

\texttt{\_apply\_time\_schedule()} 对 \texttt{torch.linspace(0,1,num\_steps+1)} 得到的时间步做 shift：

\begin{lstlisting}[language=Python]
sigma = 1 - t
shift = timestep_shift
sigma = shift * sigma / (1 + (shift - 1) * sigma)
return 1 - sigma
\end{lstlisting}

所以源码里所谓 \texttt{timestep\_shift}，不是换一个采样器，而是重新映射 \(t\) 的分布。

\subsubsection{3. T2I 文本输入和图像位置索引：源码行 431--466}

\texttt{\_build\_t2i\_query()} 先用 conversation template 拼 prompt：

\begin{lstlisting}[language=Python]
template = get_conv_template(self.template)
template.system_message = self.system_message if system_message is None else system_message
template.append_message(template.roles[0], prompt_text)
template.append_message(template.roles[1], None)
return template.get_prompt() + append_text
\end{lstlisting}

\texttt{\_build\_t2i\_text\_inputs()} 把纯文本 query token 化，并构造 \texttt{t/h/w} 三轴索引：

\begin{lstlisting}[language=Python]
t_idx = torch.arange(0, input_ids.shape[1], dtype=torch.long, device=input_ids.device)
h_idx = torch.zeros_like(t_idx)
w_idx = torch.zeros_like(t_idx)
indexes = torch.stack([t_idx, h_idx, w_idx], dim=0)
attention_mask = {"full_attention": create_block_causal_mask(indexes[0])}
\end{lstlisting}

纯文本没有图像二维坐标，所以 \texttt{h/w} 全是 0。

\texttt{\_build\_t2i\_image\_indexes()} 则给将要生成的图像 token 建位置：

\begin{lstlisting}[language=Python]
t_image = torch.full((token_h * token_w,), text_len, dtype=torch.long, device=device)
idx = torch.arange(token_h * token_w, device=device, dtype=torch.long)
h_image = idx // token_w
w_image = idx % token_w
return torch.stack([t_image, h_image, w_image], dim=0)
\end{lstlisting}

这一段说明生成图像 token 仍然进入 Qwen3 的 \texttt{t/h/w} 位置体系：同一张图里的 token 拥有同一个 \texttt{t}，但有不同的 \texttt{h/w}。

\subsubsection{4. Prefix forward：源码行 459--476}

T2I 前缀用 \texttt{input\_ids}：

\begin{lstlisting}[language=Python]
out = self.language_model.model(
    input_ids=input_ids,
    indexes=indexes,
    attention_mask=attention_mask,
    use_cache=True,
)
return out.past_key_values, out.last_hidden_state
\end{lstlisting}

IT2I / interleave 前缀用 \texttt{inputs\_embeds}，因为其中可能已经塞入视觉 embedding：

\begin{lstlisting}[language=Python]
out = self.language_model.model(
    inputs_embeds=input_imbeds,
    indexes=indexes,
    attention_mask=attention_mask,
    use_cache=True,
    image_gen_indicators=gen_indicators.view(1, -1) if gen_indicators is not None else None
)
\end{lstlisting}

这就是源码里 T2I 和 editing 的关键分叉：\key{T2I prefix 是纯文本 token；editing prefix 是已经替换过 \texttt{<IMG\_CONTEXT>} 的 embedding 序列。}

\subsubsection{5. \texttt{\_generate\_think()}：源码行 507--560}

如果打开 \texttt{think\_mode}，\texttt{t2i\_generate()} / \texttt{it2i\_generate()} 不会马上进图像 loop，而是先从 prefix 的 logits 开始自回归生成 \texttt{<think>...</think>}：

\begin{lstlisting}[language=Python]
next_token = torch.argmax(prefix_outputs.logits[:, -1, :], dim=-1)
for _ in range(max_think_tokens):
    token_item = next_token.item()
    if token_item == eos_token_id:
        break
    if token_item == think_end_token_id:
        ... append </think> to cache ...
        break
    ... append generated token to cache ...
\end{lstlisting}

生成完 think block 后，源码再把 \texttt{"\textbackslash n\textbackslash n" + IMG\_START\_TOKEN} 追加进 cache：

\begin{lstlisting}[language=Python]
append_ids = tokenizer('\n\n' + IMG_START_TOKEN, return_tensors='pt', add_special_tokens=False)['input_ids']
t_idx = self._append_text_tokens_to_cache(past_key_values, t_idx, append_ids)
\end{lstlisting}

所以 \texttt{think\_mode} 的源码含义是：\textbf{先让模型生成思考文本，并把这段思考也写进 prefix cache，再开始图像生成。}

\subsubsection{6. \texttt{\_t2i\_predict\_v()}：源码行 562--602}

这是所有图像采样 loop 复用的核心函数。输入是当前图像 token embedding、图像位置索引、prefix cache、当前时间 \(t\)、当前 patch 状态 \texttt{z}：

\begin{lstlisting}[language=Python]
outputs = self.language_model.model(
    inputs_embeds=input_embeds,
    image_gen_indicators=torch.ones(..., dtype=torch.bool),
    indexes=indexes_image,
    attention_mask=attn_mask,
    past_key_values=past_key_values,
    update_cache=False,
    use_cache=True,
)
\end{lstlisting}

这里 \texttt{update\_cache=False} 很关键：每个采样步只是临时把当前图像 token 接到 prefix 后预测速度，不把每一步的噪声图永久追加到 cache。

接着 \texttt{fm\_head} 把图像 token hidden state 投成 \texttt{x\_pred}：

\begin{lstlisting}[language=Python]
x_pred = self.fm_modules["fm_head"](
    outputs.last_hidden_state[:, -image_token_num:].view(B, L, -1)
).view(B, L, -1)
\end{lstlisting}

最后得到 velocity：

\begin{lstlisting}[language=Python]
v_pred = (x_pred - z) / (1 - t).clamp_min(self.config.t_eps)
return v_pred
\end{lstlisting}

到这里为止，源码已经告诉我们图像生成 loop 的最小单元：\texttt{image\_embeds + prefix cache -> Qwen3 -> fm\_head -> v\_pred}。

\subsubsection{7. \texttt{\_build\_it2i\_inputs()}：源码行 604--623}

editing / interleave 的输入图像不是直接作为图片对象送进 Qwen3，而是先展开成 \texttt{<IMG\_CONTEXT>} 槽位，再把视觉 embedding 写进去：

\begin{lstlisting}[language=Python]
indexes = self.get_thw_indexes(input_ids[0], grid_hw)
attention_mask = {"full_attention": create_block_causal_mask(indexes[0])}
input_embeds = self.language_model.get_input_embeddings()(input_ids)

vit_embeds = self.extract_feature(pixel_values, grid_hw=grid_hw)
selected = (input_ids == self.img_context_token_id)
input_embeds[selected] = vit_embeds.reshape(-1, C).to(input_embeds.device)
\end{lstlisting}

这段和上一节 VQA pathway 是同一个机制。区别只在后续目标：VQA 后面生成文本，editing / interleave 后面可能继续进入图像采样 loop。

\subsection{四、\texttt{interleave\_gen():954--1295}：源码先写 interleave，所以先按它读}

\subsubsection{1. 行 979--1044：准备 special token、输入图像和 condition cache}

函数一开始设置 special token id 和 \texttt{t\_eps}：

\begin{lstlisting}[language=Python]
self.img_context_token_id = tokenizer.convert_tokens_to_ids(IMG_CONTEXT_TOKEN)
self.img_start_token_id = tokenizer.convert_tokens_to_ids(IMG_START_TOKEN)
self.config.t_eps = t_eps
\end{lstlisting}

如果传了输入图像，源码先补足 prompt 里的 \texttt{<image>}，再逐张调用 \texttt{load\_image\_native()}：

\begin{lstlisting}[language=Python]
image_token_count = prompt.count('<image>')
assert len(images) >= image_token_count
if len(images) > image_token_count:
    prompt = "<image>\n" * (len(images) - image_token_count) + prompt

for image in images:
    cur_pixel_values, cur_grid_hw = load_image_native(...)
\end{lstlisting}

接着构造 condition query，并把每个 \texttt{<image>} 替换成 \texttt{<img>} + 多个 \texttt{<IMG\_CONTEXT>} + \texttt{</img>}：

\begin{lstlisting}[language=Python]
num_patch_token = int(grid_hw_list[i][0, 0] * grid_hw_list[i][0, 1] * self.downsample_ratio**2)
image_tokens = IMG_START_TOKEN + IMG_CONTEXT_TOKEN * num_patch_token + IMG_END_TOKEN
query = query.replace('<image>', image_tokens, 1)
\end{lstlisting}

然后 \texttt{\_build\_it2i\_inputs()} 把输入图片视觉 embedding 写入槽位，再跑一次 LLM 得到 condition cache：

\begin{lstlisting}[language=Python]
input_embeds_cond, indexes_cond, attention_mask_cond = self._build_it2i_inputs(...)
outputs_cond = self.language_model(inputs_embeds=input_embeds_cond, indexes=indexes_cond, attention_mask=attention_mask_cond, use_cache=True)
past_key_values_cond = outputs_cond.past_key_values
\end{lstlisting}

\subsubsection{2. 行 1046--1065：准备 text-uncondition 和 image-uncondition cache}

interleave 图像生成时有三条 cache：

\begin{lstlisting}[language=Python]
# text uncondition: only keep input images
question_text_uncondition = '<image>' * len(images)
...
past_key_values_tu = outputs_tu.past_key_values

# image uncondition: empty T2I query
query_img_uncond = self._build_t2i_query("", append_text=IMG_START_TOKEN)
...
past_key_values_iu = outputs_iu.past_key_values
\end{lstlisting}

这里不要把它理解成抽象概念，按源码看就是三份 \texttt{past\_key\_values}：

\begin{itemize}[leftmargin=1.5em]
  \item \textbf{\texttt{past\_key\_values\_cond}：} 当前完整上下文；
  \item \textbf{\texttt{past\_key\_values\_tu}：} 去掉文本条件，只保留输入图像；
  \item \textbf{\texttt{past\_key\_values\_iu}：} 空图像生成条件。
\end{itemize}

\subsubsection{3. 行 1072--1127：文本自回归，直到遇到 \texttt{<img>}}

主循环先从 \texttt{outputs\_cond.logits} 取下一个 token：

\begin{lstlisting}[language=Python]
next_token = torch.argmax(outputs_cond.logits[:, -1, :], dim=-1)
\end{lstlisting}

只要下一个 token 不是 EOS，也不是 \texttt{<img>}，就按普通 LLM decoding 追加到 condition cache：

\begin{lstlisting}[language=Python]
while True:
    token_item = next_token.item()
    if token_item == eos_token_id or token_item == self.img_start_token_id:
        break
    ...
    outputs_cond = self.language_model(
        input_ids=next_token.unsqueeze(0),
        past_key_values=past_key_values_cond,
        use_cache=True
    )
\end{lstlisting}

遇到 \texttt{<img>} 后，源码先把这个 token 写入 condition 和 text-uncondition 两个 cache：

\begin{lstlisting}[language=Python]
outputs_cond = self.language_model(input_ids=next_token.unsqueeze(0), past_key_values=past_key_values_cond, use_cache=True)
outputs_tu = self.language_model(input_ids=next_token.unsqueeze(0), past_key_values=past_key_values_tu, use_cache=True)
\end{lstlisting}

到这一行，控制流才真正从“文本生成”切到“图像生成”。

\subsubsection{4. 行 1128--1178：为即将生成的一张图准备位置、噪声和 KV cache}

源码根据当前第几张图选择 \texttt{image\_size}，并计算生成图像 token 网格：

\begin{lstlisting}[language=Python]
token_h = image_size[1] // (self.patch_size * merge_size)
token_w = image_size[0] // (self.patch_size * merge_size)
indexes_image_condition = self._build_t2i_image_indexes(token_h, token_w, t_index_cond + 1, device=device)
\end{lstlisting}

然后按分辨率初始化随机噪声图：

\begin{lstlisting}[language=Python]
noise_scale = self.noise_scale
if self.noise_scale_mode in ("resolution", "dynamic", 'dynamic_sqrt'):
    noise_scale = math.sqrt((grid_h*grid_w)/(merge_size**2)/base) * float(self.noise_scale)
image_prediction = noise_scale * torch.randn((1, 3, image_size[1], image_size[0]), device=device, dtype=outputs_cond.logits.dtype, generator=generator)
\end{lstlisting}

\texttt{prepare\_flash\_kv\_cache()} 提前给三条 cache 预留当前图像 token 长度：

\begin{lstlisting}[language=Python]
prepare_flash_kv_cache(past_key_values_cond_cfg, current_len=token_h * token_w, batch_size=1)
prepare_flash_kv_cache(past_key_values_tu_cfg, current_len=token_h * token_w, batch_size=1)
prepare_flash_kv_cache(past_key_values_iu_cfg, current_len=token_h * token_w, batch_size=1)
\end{lstlisting}

\subsubsection{5. 行 1179--1228：图像 flow loop}

每个 step 的源码顺序固定：先取 \(t,t_{next}\)，再从当前图像构造 \texttt{z} 和 \texttt{image\_embeds}：

\begin{lstlisting}[language=Python]
z = self.patchify(image_prediction, self.patch_size * merge_size)
image_input = self.patchify(image_prediction, self.patch_size, channel_first=True)
image_embeds = self.extract_feature(
    image_input.view(1 * grid_h*grid_w, -1),
    gen_model=True,
    grid_hw=gen_grid_hw,
).view(1, token_h*token_w, -1)
\end{lstlisting}

再加时间 embedding 和可选噪声强度 embedding：

\begin{lstlisting}[language=Python]
timestep_embeddings = self.fm_modules['timestep_embedder'](t_expanded).view(1, token_h*token_w, -1)
if self.add_noise_scale_embedding:
    timestep_embeddings += noise_embeddings
image_embeds = image_embeds + timestep_embeddings
\end{lstlisting}

接着调用 \texttt{\_t2i\_predict\_v()}。根据 \texttt{cfg\_scale} / \texttt{img\_cfg\_scale}，源码可能跑一条、两条或三条 cache：

\begin{lstlisting}[language=Python]
out_cond = self._t2i_predict_v(... past_key_values_cond_cfg ...)
...
v_pred = (
    out_uncond
    + cfg_scale * (out_cond - out_img_cond)
    + img_cfg_scale * (out_img_cond - out_uncond)
)
\end{lstlisting}

最后直接 Euler 更新并拼回图像：

\begin{lstlisting}[language=Python]
z = z + (t_next - t) * v_pred
image_prediction = self.unpatchify(z, self.patch_size * merge_size, image_size[1], image_size[0])
\end{lstlisting}

\subsubsection{6. 行 1237--1293：生成图像重新编码回文本上下文}

interleave 和 T2I / editing 最大的源码差异在这里。图像生成完并不是直接返回，而是先转回 understanding branch 的归一化格式：

\begin{lstlisting}[language=Python]
pred_img = image_prediction[0].unsqueeze(0).to(torch.bfloat16)
raw_img = pred_img * 0.5 + 0.5
und_img = (raw_img - img_mean) / img_std
\end{lstlisting}

然后重新切 patch、抽 understanding 视觉特征：

\begin{lstlisting}[language=Python]
flatten_pixel_values = ...
vit_embeds = self.extract_feature(flatten_pixel_values, grid_hw=gen_grid_hw[:1]).unsqueeze(0)
\end{lstlisting}

最后把这张新图的视觉 token 和 \texttt{</img>} token 追加进 cache：

\begin{lstlisting}[language=Python]
inputs_embeds_img = torch.cat([vit_embeds, img_end_embed], dim=1)
outputs = self.language_model(
    inputs_embeds=inputs_embeds_img,
    indexes=indexes,
    attention_mask=attention_mask_dict,
    past_key_values=cache,
    use_cache=True
)
\end{lstlisting}

源码层面的含义很明确：\textbf{interleave 生成一张图以后，会马上把这张图重新当作输入图像理解一遍，使后续文本生成可以看见它。}

\subsection{五、\texttt{it2i\_generate():1298--1616}：editing 按源码块读}

\subsubsection{1. 行 1301--1328：special token、补 \texttt{<image>}、加载输入图像}

editing 函数先设置 \texttt{img\_context\_token\_id} 和 \texttt{t\_eps}：

\begin{lstlisting}[language=Python]
self.img_context_token_id = tokenizer.convert_tokens_to_ids(IMG_CONTEXT_TOKEN)
self.config.t_eps = t_eps
\end{lstlisting}

然后检查 prompt 里有几个 \texttt{<image>}。如果传入图片更多，源码会自动补占位符：

\begin{lstlisting}[language=Python]
image_token_count = prompt.count('<image>')
assert len(images) >= image_token_count
if len(images) > image_token_count:
    if image_token_count == 0 and len(images) > 1:
        prompt = "".join(f"Image-{i + 1}:<image>\n" for i in range(len(images))) + prompt
    else:
        prompt = "<image>\n" * (len(images) - image_token_count) + prompt
\end{lstlisting}

每张输入图像走 \texttt{load\_image\_native()}，得到 \texttt{pixel\_values} 和 \texttt{grid\_hw}：

\begin{lstlisting}[language=Python]
cur_pixel_values, cur_grid_hw = load_image_native(
    image,
    self.patch_size,
    self.downsample_ratio,
    min_pixels=512 * 512,
    max_pixels=min(2048*2048, (4096 * 4096) // len(images)),
    upscale=False,
)
\end{lstlisting}

\subsubsection{2. 行 1330--1368：构造三种 query，并把输入图像写入 embedding}

editing 的 CFG 分支在源码中先由三个布尔量决定：

\begin{lstlisting}[language=Python]
needs_cfg = not (cfg_scale == 1 and img_cfg_scale == 1)
needs_img_condition = needs_cfg and (img_cfg_scale == 1 or cfg_scale != img_cfg_scale)
needs_uncondition = needs_cfg and img_cfg_scale != 1
\end{lstlisting}

接着生成三种 query：

\begin{lstlisting}[language=Python]
query_condition = self._build_t2i_query(question_condition, system_message=SYSTEM_MESSAGE_FOR_GEN, append_text=think_content)
query_img_condition = self._build_t2i_query('<image>' * len(images), append_text=IMG_START_TOKEN) if needs_img_condition else None
query_uncondition = self._build_t2i_query("", append_text=IMG_START_TOKEN) if needs_uncondition else None
\end{lstlisting}

然后把 \texttt{<image>} 替换成图像 token 槽位：

\begin{lstlisting}[language=Python]
num_patch_token = int(grid_hw[i, 0] * grid_hw[i, 1] * self.downsample_ratio**2)
image_tokens = IMG_START_TOKEN + IMG_CONTEXT_TOKEN * num_patch_token + IMG_END_TOKEN
query_condition = query_condition.replace('<image>', image_tokens, 1)
\end{lstlisting}

最后调用 \texttt{\_build\_it2i\_inputs()}，把输入图像 embedding 写进这些槽位：

\begin{lstlisting}[language=Python]
input_embeds_condition, indexes_condition, attention_mask_condition_prefix = self._build_it2i_inputs(
    tokenizer, query_condition, pixel_values, grid_hw
)
\end{lstlisting}

\subsubsection{3. 行 1369--1425：生成图像位置索引，并跑 prefix forward}

输出图尺寸决定目标图像 token 数：

\begin{lstlisting}[language=Python]
token_h = image_size[1] // (self.patch_size * merge_size)
token_w = image_size[0] // (self.patch_size * merge_size)
\end{lstlisting}

每个分支都有自己的图像位置索引，因为每个 prefix 的 \texttt{t} 长度不同：

\begin{lstlisting}[language=Python]
indexes_image_condition = self._build_t2i_image_indexes(
    token_h, token_w, indexes_condition[0].max() + 1, device=input_embeds_condition.device
)
\end{lstlisting}

如果 \texttt{think\_mode=True}，condition 分支先生成 think，再重新计算图像位置索引；否则直接用 \texttt{\_it2i\_prefix\_forward()} 建 cache：

\begin{lstlisting}[language=Python]
past_key_values_condition, hidden_states_condition = self._it2i_prefix_forward(
    input_embeds_condition, indexes_condition, attention_mask_condition_prefix
)
\end{lstlisting}

image-only 和 unconditional 分支也分别跑 prefix forward：

\begin{lstlisting}[language=Python]
past_key_values_img_condition, _ = self._it2i_prefix_forward(...)
past_key_values_uncondition, _ = self._it2i_prefix_forward(...)
\end{lstlisting}

\subsubsection{4. 行 1437--1500：batch 扩展 cache、准备 flash cache、初始化噪声图}

源码先把每层 KV cache expand 到 \texttt{batch\_size}：

\begin{lstlisting}[language=Python]
past_key_values_condition.layers[layer_idx].keys = past_key_values_condition.layers[layer_idx].keys.expand(
    batch_size, *past_key_values_condition.layers[layer_idx].keys.shape[1:]
)
\end{lstlisting}

再调用 \texttt{prepare\_flash\_kv\_cache()} 给将要生成的图像 token 预留空间：

\begin{lstlisting}[language=Python]
prepare_flash_kv_cache(
    past_key_values_condition,
    current_len=token_h * token_w,
    batch_size=batch_size,
)
\end{lstlisting}

然后根据输出分辨率初始化当前图像状态：

\begin{lstlisting}[language=Python]
grid_h = image_size[1] // self.patch_size
grid_w = image_size[0] // self.patch_size
grid_hw = torch.tensor([[grid_h, grid_w]] * batch_size, device=device)

image_prediction = noise_scale * torch.randn(
    (batch_size, 3, image_size[1], image_size[0]), device=device, dtype=dtype, generator=generator
)
\end{lstlisting}

\subsubsection{5. 行 1502--1605：editing 的采样 loop}

采样 loop 开头和 interleave 相同：当前图像先切成 \texttt{z} 和 vision encoder 输入：

\begin{lstlisting}[language=Python]
z = self.patchify(image_prediction, self.patch_size * merge_size)
image_input = self.patchify(image_prediction, self.patch_size, channel_first=True)
image_embeds = self.extract_feature(
    image_input.view(batch_size * grid_h * grid_w, -1),
    gen_model=True,
    grid_hw=grid_hw,
).view(batch_size, token_h * token_w, -1)
\end{lstlisting}

然后加时间条件：

\begin{lstlisting}[language=Python]
timestep_embeddings = self.fm_modules['timestep_embedder'](t_expanded).view(batch_size, token_h * token_w, -1)
image_embeds = image_embeds + timestep_embeddings
\end{lstlisting}

condition 分支总会先跑：

\begin{lstlisting}[language=Python]
out_cond = self._t2i_predict_v(
    image_embeds,
    indexes_image_condition,
    attention_mask_condition,
    past_key_values_condition,
    t,
    z,
    image_token_num=token_h * token_w,
)
\end{lstlisting}

之后源码按 \texttt{cfg\_scale} / \texttt{img\_cfg\_scale} 分情况组合：

\begin{lstlisting}[language=Python]
if not use_cfg:
    v_pred = out_cond
elif img_cfg_scale == 1:
    v_pred = out_img_cond + cfg_scale * (out_cond - out_img_cond)
elif cfg_scale == img_cfg_scale:
    v_pred = out_uncond + cfg_scale * (out_cond - out_uncond)
else:
    v_pred = (
        out_uncond
        + cfg_scale * (out_cond - out_img_cond)
        + img_cfg_scale * (out_img_cond - out_uncond)
    )
\end{lstlisting}

最后是同一行 Euler 更新和反 patch：

\begin{lstlisting}[language=Python]
z = z + (t_next - t) * v_pred
image_prediction = self.unpatchify(z, self.patch_size * merge_size, image_size[1], image_size[0])
\end{lstlisting}

\subsubsection{6. 行 1607--1616：清理 cache 并返回}

loop 结束后，源码清理 flash cache，并根据 \texttt{think\_mode} 决定返回一张图，还是图像加 think 文本：

\begin{lstlisting}[language=Python]
clear_flash_kv_cache(past_key_values_condition)
...
self.last_think_content = think_text
if think_mode:
    return image_prediction, think_text
return image_prediction
\end{lstlisting}

\subsection{六、\texttt{t2i\_generate():1619--1787}：T2I 按源码块读}

\subsubsection{1. 行 1620--1640：构造 condition / uncondition 文本 query}

T2I 比 editing 少了输入图片，所以一开始只有文本：

\begin{lstlisting}[language=Python]
question_condition = f"{prompt}"
needs_cfg = cfg_scale > 1
think_content = '<think>\n' if think_mode else '<think>\n\n</think>\n\n' + IMG_START_TOKEN
query_condition = self._build_t2i_query(question_condition, system_message=SYSTEM_MESSAGE_FOR_GEN, append_text=think_content)
query_uncondition = self._build_t2i_query("", append_text=IMG_START_TOKEN) if needs_cfg else None
\end{lstlisting}

再把 query 变成 \texttt{input\_ids/indexes/attention\_mask}：

\begin{lstlisting}[language=Python]
input_ids_condition, indexes_condition, attention_mask_condition_prefix = self._build_t2i_text_inputs(tokenizer, query_condition)
\end{lstlisting}

这和 editing 的差异非常具体：T2I 走 \texttt{\_build\_t2i\_text\_inputs()}，editing 走 \texttt{\_build\_it2i\_inputs()}。

\subsubsection{2. 行 1642--1678：目标图像位置索引和 prefix cache}

T2I 根据输出尺寸计算目标图 token 网格：

\begin{lstlisting}[language=Python]
token_h = image_size[1] // (self.patch_size * merge_size)
token_w = image_size[0] // (self.patch_size * merge_size)
\end{lstlisting}

再构造 condition / uncondition 的目标图像位置索引：

\begin{lstlisting}[language=Python]
indexes_image_condition = self._build_t2i_image_indexes(token_h, token_w, indexes_condition.shape[1], device=input_ids_condition.device)
\end{lstlisting}

如果 \texttt{think\_mode=False}，直接跑 prefix forward：

\begin{lstlisting}[language=Python]
past_key_values_condition, hidden_states_condition = self._t2i_prefix_forward(
    input_ids_condition,
    indexes_condition,
    attention_mask_condition_prefix,
)
\end{lstlisting}

如果有 CFG，无条件 query 也跑一次 prefix forward：

\begin{lstlisting}[language=Python]
past_key_values_uncondition, _ = self._t2i_prefix_forward(
    input_ids_uncondition,
    indexes_uncondition,
    attention_mask_uncondition_prefix,
)
\end{lstlisting}

\subsubsection{3. 行 1687--1728：扩展 cache、准备随机图像初值和时间步}

batch 维度是通过 expand cache 实现的：

\begin{lstlisting}[language=Python]
past_key_values_condition.layers[layer_idx].keys = past_key_values_condition.layers[layer_idx].keys.expand(
    batch_size, *past_key_values_condition.layers[layer_idx].keys.shape[1:]
)
\end{lstlisting}

接着预留 flash KV cache：

\begin{lstlisting}[language=Python]
prepare_flash_kv_cache(
    past_key_values_condition,
    current_len=token_h * token_w,
    batch_size=batch_size,
)
\end{lstlisting}

然后初始化噪声图：

\begin{lstlisting}[language=Python]
grid_h = image_size[1] // self.patch_size
grid_w = image_size[0] // self.patch_size
grid_hw = torch.tensor([[grid_h, grid_w]]*batch_size, device=device)

image_prediction = noise_scale * torch.randn(
    (batch_size, 3, image_size[1], image_size[0]),
    device=device,
    dtype=dtype,
    generator=generator,
)
\end{lstlisting}

最后创建时间表：

\begin{lstlisting}[language=Python]
timesteps = torch.linspace(0.0, 1.0, num_steps+1, device=device)
if enable_timestep_shift:
    timesteps = self._apply_time_schedule(timesteps, token_h*token_w, timestep_shift)
\end{lstlisting}

\subsubsection{4. 行 1730--1745：T2I loop 前半段，当前图像变成 token embedding}

每一步先把当前图像拆成两份：

\begin{lstlisting}[language=Python]
z = self.patchify(image_prediction, self.patch_size * merge_size)
image_input = self.patchify(image_prediction, self.patch_size, channel_first=True)
\end{lstlisting}

\texttt{z} 是要被更新的 RGB patch 状态；\texttt{image\_input} 送给生成侧 vision encoder：

\begin{lstlisting}[language=Python]
image_embeds = self.extract_feature(
    image_input.view(batch_size * grid_h*grid_w, -1),
    gen_model=True,
    grid_hw=grid_hw,
).view(batch_size, token_h*token_w, -1)
\end{lstlisting}

时间步 embedding 直接加到图像 token 上：

\begin{lstlisting}[language=Python]
timestep_embeddings = self.fm_modules['timestep_embedder'](t_expanded).view(batch_size, token_h*token_w, -1)
image_embeds = image_embeds + timestep_embeddings
\end{lstlisting}

然后 condition 分支预测速度：

\begin{lstlisting}[language=Python]
v_pred_condition = self._t2i_predict_v(
    image_embeds,
    indexes_image_condition,
    attention_mask_condition,
    past_key_values_condition,
    t,
    z,
    image_token_num=token_h*token_w,
)
\end{lstlisting}

\subsubsection{5. 行 1747--1774：T2I CFG 分支}

如果当前 \(t\) 落在 \texttt{cfg\_interval}，且 \texttt{cfg\_scale > 1}，源码会再跑 uncondition 分支：

\begin{lstlisting}[language=Python]
v_pred_uncondition = self._t2i_predict_v(
    image_embeds,
    indexes_image_uncondition,
    attention_mask_uncondition,
    past_key_values_uncondition,
    t,
    z,
    image_token_num=token_h*token_w,
)
\end{lstlisting}

普通 CFG 是：

\begin{lstlisting}[language=Python]
v_pred = v_pred_uncondition + cfg_scale * (v_pred_condition - v_pred_uncondition)
\end{lstlisting}

如果 \texttt{cfg\_norm == 'global'} 或 \texttt{'channel'}，源码会把 CFG 后的速度范数往 condition 速度范数上拉：

\begin{lstlisting}[language=Python]
norm_v_condition = torch.norm(v_pred_condition, dim=(1,2), keepdim=True)
norm_v_cfg = torch.norm(v_pred, dim=(1,2), keepdim=True)
scale = (norm_v_condition / (norm_v_cfg + 1e-8)).clamp(min=0, max=1.0)
v_pred = v_pred * scale
\end{lstlisting}

如果不满足 CFG 条件，就直接用 condition 速度：

\begin{lstlisting}[language=Python]
v_pred = v_pred_condition
\end{lstlisting}

\subsubsection{6. 行 1776--1787：Euler 更新、反 patch、返回}

T2I 最后一段和 editing 相同：

\begin{lstlisting}[language=Python]
z = z + (t_next - t) * v_pred
image_prediction = self.unpatchify(z, self.patch_size * merge_size, image_size[1], image_size[0])
\end{lstlisting}

loop 结束后清理 cache 并返回：

\begin{lstlisting}[language=Python]
clear_flash_kv_cache(past_key_values_condition)
if past_key_values_uncondition is not None:
    clear_flash_kv_cache(past_key_values_uncondition)

if think_mode:
    return image_prediction, think_text
return image_prediction
\end{lstlisting}

\subsection{七、按源码顺序串一次调用栈}

\subsubsection{1. T2I}

\begin{lstlisting}[language=Python]
examples/t2i/inference.py
  SenseNovaU1T2I.generate()
    -> model.t2i_generate()                         # modeling_neo_chat.py:1619
       -> _build_t2i_query()
       -> _build_t2i_text_inputs()
       -> _build_t2i_image_indexes()
       -> _t2i_prefix_forward()
       -> optional _generate_think()
       -> initialize image_prediction from randn
       -> for step in timesteps:
            patchify(image_prediction)              # z
            patchify(image_prediction, channel_first=True)
            extract_feature(..., gen_model=True)
            timestep_embedder(t)
            _t2i_predict_v(... condition cache ...)
            optional _t2i_predict_v(... uncondition cache ...)
            CFG on v_pred
            z = z + (t_next - t) * v_pred
            unpatchify(z)
\end{lstlisting}

\subsubsection{2. Editing}

\begin{lstlisting}[language=Python]
examples/editing/inference.py
  SenseNovaU1Editing.edit()
    -> model.it2i_generate()                        # modeling_neo_chat.py:1298
       -> load_image_native(input images)
       -> replace <image> with <IMG_CONTEXT> slots
       -> _build_it2i_inputs()
            -> extract_feature(pixel_values)        # understanding branch
            -> input_embeds[selected] = vit_embeds
       -> _it2i_prefix_forward()                    # condition / img_condition / uncondition
       -> initialize image_prediction from randn
       -> for step in timesteps:
            same image flow loop as T2I
            combine out_cond / out_img_cond / out_uncond
            unpatchify(z)
\end{lstlisting}

\subsubsection{3. Interleave}

\begin{lstlisting}[language=Python]
examples/interleave/inference.py
  SenseNovaU1Interleave.generate()
    -> model.interleave_gen()                       # modeling_neo_chat.py:954
       -> build condition / text-uncondition / image-uncondition caches
       -> while True:
            generate text tokens with condition cache
            if next token is <img>:
                append <img> to caches
                run image flow loop
                re-encode generated image with understanding encoder
                append generated image tokens back into cache
            if eos: break
\end{lstlisting}

\subsection{八、本节小结：只保留源码能直接支持的结论}

\begin{summarybox}
按源码读下来，可以得到四个确定结论：\\
1. 三个 example wrapper 都很薄，分别进入 \texttt{t2i\_generate()}、\texttt{it2i\_generate()}、\texttt{interleave\_gen()}；\\
2. 图像采样状态是 \texttt{image\_prediction}，通过 \texttt{patchify()} 变成 RGB patch 序列 \texttt{z}，再通过 \texttt{unpatchify()} 回到图像；\\
3. 每个采样 step 都会调用 \texttt{extract\_feature(..., gen\_model=True)}、\texttt{timestep\_embedder} 和 \texttt{\_t2i\_predict\_v()}，其中 \texttt{\_t2i\_predict\_v()} 负责用 Qwen3 + \texttt{fm\_head} 得到 \texttt{v\_pred = (x\_pred - z)/(1-t)}；\\
4. T2I、editing、interleave 的差异不是采样 loop 本身，而是源码里进入 loop 前准备了哪些 \texttt{past\_key\_values}，以及 loop 后是否像 interleave 那样把生成图像重新编码回 cache。
\end{summarybox}

% ============================================================
\subsection{参考资料}
% ============================================================

\begingroup
\small
\sloppy
\begin{itemize}[leftmargin=1.5em]
  \item OpenSenseNova，\emph{SenseNova-U1 官方 GitHub 仓库}。\\
  \url{https://github.com/OpenSenseNova/SenseNova-U1}
  \item OpenSenseNova，\emph{SenseNova-U1 examples README：T2I、editing、interleaved generation 运行参数说明}。\\
  \url{https://github.com/OpenSenseNova/SenseNova-U1/blob/main/examples/README.md}
  \item OpenSenseNova，\emph{Text-to-Image inference wrapper：\texttt{SenseNovaU1T2I.generate}}。\\
  \url{https://github.com/OpenSenseNova/SenseNova-U1/blob/main/examples/t2i/inference.py}
  \item OpenSenseNova，\emph{Image editing inference wrapper：\texttt{SenseNovaU1Editing.edit}}。\\
  \url{https://github.com/OpenSenseNova/SenseNova-U1/blob/main/examples/editing/inference.py}
  \item OpenSenseNova，\emph{Interleaved generation wrapper：\texttt{SenseNovaU1Interleave.generate}}。\\
  \url{https://github.com/OpenSenseNova/SenseNova-U1/blob/main/examples/interleave/inference.py}
  \item OpenSenseNova，\emph{NEOChatModel generation methods：\texttt{patchify}、\texttt{unpatchify}、\texttt{interleave\_gen}、\texttt{it2i\_generate}、\texttt{t2i\_generate}、\texttt{\_t2i\_predict\_v}}。\\
  \url{https://github.com/OpenSenseNova/SenseNova-U1/blob/main/src/sensenova_u1/models/neo_unify/modeling_neo_chat.py}
  \item OpenSenseNova，\emph{Flow-matching helper modules：\texttt{TimestepEmbedder}、\texttt{FlowMatchingHead}、\texttt{ConvDecoder}}。\\
  \url{https://github.com/OpenSenseNova/SenseNova-U1/blob/main/src/sensenova_u1/models/neo_unify/modeling_fm_modules.py}
  \item Yaron Lipman et al.，\emph{Flow Matching for Generative Modeling}。\\
  \url{https://arxiv.org/abs/2210.02747}
  \item Xingchao Liu et al.，\emph{Flow Straight and Fast: Learning to Generate and Transfer Data with Rectified Flow}。\\
  \url{https://arxiv.org/abs/2209.03003}
  \item Jonathan Ho and Tim Salimans，\emph{Classifier-Free Diffusion Guidance}。\\
  \url{https://arxiv.org/abs/2207.12598}
  \item William Peebles and Saining Xie，\emph{Scalable Diffusion Models with Transformers}。\\
  \url{https://arxiv.org/abs/2212.09748}
\end{itemize}
\endgroup

